
%\documentclass[oneside,final,14pt]{extreport}
\documentclass[12pt, a4paper, openany]{book}



\usepackage[left=2cm,right=2cm,top=2cm,bottom=2cm,bindingoffset=0cm]{geometry}
%\usepackage[koi8-r]{inputenc}
%\usepackage[russianb]{babel}
\usepackage{vmargin}

\setpapersize{A4}
\usepackage[T2A]{fontenc}
\usepackage[utf8x]{inputenc}
\usepackage[english, russian]{babel}
\setmarginsrb{2cm}{2cm}{2cm}{2cm}{0pt}{5mm}{0pt}{0mm}
\usepackage{indentfirst}
\usepackage{nicefrac} % For comparison
%\usepackage{xfrac} % Works better with other fonts
%\usepackage[unicode, pdftex]{hyperref}
\usepackage{lettrine}
\usepackage[usenames]{color}
\usepackage{colortbl}
\usepackage{mathtext}
\usepackage{epigraph}
\usepackage{amsmath, amsfonts, amssymb, mathrsfs}
%\usepackage{mathptmx}
%\usepackage{txfonts}
\usepackage{pxfonts}
\usepackage[pagestyles]{titlesec}
\usepackage{ebgaramond}
\usepackage{awesomebox}
\usepackage{enumitem}
\usepackage{makeidx}
\usepackage[letterspace=150]{microtype}
\makeindex
\usepackage{etoolbox}
\makeatletter
\newlength\epitextskip
\pretocmd{\@epitext}{\em}{}{}
\apptocmd{\@epitext}{\em}{}{}
\patchcmd{\epigraph}{\@epitext{#1}\\}{\@epitext{#1}\\[\epitextskip]}{}{}
\makeatother
%running fraction with slash - requires math mode.
\newcommand*\rfrac[2]{{}^{#1}\!/_{#2}}

\DeclareSymbolFont{Xlargesymbols}{OMX}{cmex}{m}{n}
\DeclareMathSymbol{\Xsum}{\mathop}{Xlargesymbols}{80}

\setlength\epigraphrule{0pt}
\setlength\epitextskip{2ex}
\setlength\epigraphwidth{.8\textwidth}

\usepackage{xfrac} % Works better with other fonts
\usepackage[colorlinks=true,linkcolor=black,urlcolor=black,bookmarksopen=true]{hyperref}

\usepackage{fancyhdr} % пакет для установки колонтитулов
\pagestyle{fancy} % смена стиля оформления страниц
\fancyhf{} % очистка текущих значений
\fancyhead[C]{\thepage} % установка верхнего колонтитула
\renewcommand{\headrulewidth}{0pt} % убрать разделительную линию


% Настройка вертикальных и горизонтальных отступов
\titlespacing{\chapter}{0pt}{5pt}{5pt}
\titlespacing{\section}{\parindent}{4mm}{4mm}
\titlespacing{\subsection}{\parindent}{3mm}{3mm}


%
%\renewcommand{\tabcolsep}{1cm} %% increase table column spacing

% Настройка задачи со зведочкой
\newcounter{namedlistcounter} % number the items
\newenvironment{withdot}
{\begin{list}
		{\arabic{namedlistcounter}*.} % labeling 
		{\usecounter{namedlistcounter} % set counter
			\setlength{\leftmargin}{3em}} % set spacing 
	}
	{\end{list}}


\newcommand{\anonsection}[1]{ \section*{#1} \addcontentsline{toc}{section}{\numberline {}#1}} 

\makeatletter %%%%% <---- Starting chapter without a pagebreak
\renewcommand\chapter{\par%
	\thispagestyle{plain}% \global\@topnum\z@
	\@afterindentfalse \secdef\@chapter\@schapter}
\makeatother %%%%% <---- Starting chapter without a pagebreak
\titleformat{\chapter}[display]
{\normalfont\bfseries}{}{0pt}{\Large}

\newpagestyle{mystyle}{
	\sethead[\thepage][][]{}{}{\thepage}	
}

\renewcommand{\rmdefault}{cmr}


\pagestyle{mystyle}
\sloppy
\begin{document}

	
	\begin{titlepage}
		
		\begin{center}
			\topskip0pt
			\vspace*{\fill}
			
			

			{\Huge\bf МАТЕМАТИКА КАК ПРОФЕССИЯ\\}
			\ \\
			\ \\
			
			{\Huge\bf (О воспитательном эффекте математического образования)\\}
			\ \\
			\ \\
			{\large\bf СБОРНИК СТАТЕЙ\\}
			
			
			
			\ \\
			
			
			\vspace*{\fill}
			
			\vfill
			
			
Издательство

«Знание»

Москва
			
1980
			

		\end{center}
		
	\end{titlepage}
	
	\thispagestyle{empty} % выключаем отображение номера для этой страницы
	
	\newpage
	

\setcounter{secnumdepth}{0}
	
	\section[\textls{Б. В. Гнеденко}. Предисловие]{{\small Б. В. ГНЕДЕНКО} \\ \textbf{ПРЕДИСЛОВИЕ}}


Каких-нибудь пятьдесят лет назад профессия математика находила в сущности лишь два направления применений —  преподавание в средней школе и научные исследования, как правило, связанные с педагогической деятельностью в университете или в высшем учебном заведении технического или же сельскохозяйственного характера. Лишь небольшая доля математиков использовала свои знания в страховых учреждениях, в вычислительных бюро и на опытных сельскохозяйственных станциях. В наши дни положение резко изменилось и, помимо только что названной традиционной работы, во всем мире для лиц с  математическим образованием открылись многочисленные новые возможности: работа в качестве исследователей в технических проектных учреждениях, в  научно-исследовательских учреждениях инженерного, медицинского, биологического, экономического и сельскохозяйственного профиля. Многие тысячи математиков работают в вычислительных центрах и непосредственно в заводских лабораториях. Математики принимают участие в расчете траектории полета космической станции и прочности сооружения при воздействии на него тех или иных нагрузок. Математики участвуют в оценке влияния погодных факторов,  способов обработки почвы на урожайность сельскохозяйственных культур. Все чаще медицинские учреждения начинают привлекать математиков в качестве постоянных сотрудников или же консультантов для решения чисто медицинских вопросов, связанных с диагностикой, организацией медицинского обслуживания, обработкой опытных данных. Математики разрабатывают методы оценки надежности технических изделий по результатам испытаний, а также качества больших партий продукции по результатам проверки сравнительно небольшой их доли. 

Труд математика теперь требуется в самых разнообразных областях деятельности. Недаром сейчас значительный реальный вес приобрело утверждение, что математика превратилась в производительную силу. И действительно своевременно и разумно примененные математические методы позволяют экономить материалы, труд, средства, находить оптимальный технологический режим. 

Однако, чтобы проводить необходимые расчеты, чтобы получить формулы или последовательность действий, по которым эти расчеты проводятся, нужно создать  математическую модель изучаемого явления, которая достаточно адекватно передавала бы его особенности. Создание математической модели явления и получение из нее логических выводов — это еще одна большая задача, стоящая перед математиком, занимающимся прикладными исследованиями. Наконец, перед математиком-прикладником стоит важная задача проверки соответствия развитой им теории изучаемому явлению. Все эти три ступени прикладной деятельности математика не приходят сами по себе. Этому искусству следует специально обучать. Можно превосходно овладеть математикой в абстрактном ее изложении и не видеть за понятиями и теоремами их прикладного значения, их реальной интерпретации. 

Нередко случается, что при решении прикладной задачи математик сталкивается с таким положением, когда нет готовых математических средств, необходимых для исследования. В этом случае приходится создавать недостающие звенья самой математической теории. 

Мы видим, таким образом, что воспитание математика-прикладника неразрывно связано с получением им основательного собственно математического образования и с воспитанием дополнительных качеств, позволяющих видеть за абстрактными формулами их реальное содержание. Среди этих свойств хотелось бы отметить увлеченность математическим познанием процессов природы, инженерных задач или же вопросов организации производства. Первым, кто закладывает основы образования математика, его математического воспитания, является учитель математики. И совершенно ясно, что его роль в настоящее время исключительно важна, но только она сделалась более многогранной по сравнению даже со сравнительно недавними временами. Преподавателю теперь, более чем когда бы то ни было ранее, необходимо заниматься не только обучением математике, но и воспитанием, в первую очередь воспитанием научного мировоззрения, понимания места математики в познании окружающего нас мира и его закономерностей. 

Вопросы обучения и воспитания учащихся (на разных уровнях обучения) во все времена интересовали ученых и педагогов и всегда были между собой тесно связаны. Выдающиеся представители математики прошлого  оставили после себя не только труды чисто математического  содержания, но и работы, относящиеся к вопросам  обучения и воспитания научного мировоззрения, трудолюбия, высоких моральных качеств, чувства ответственности перед страной и народом. Темы образования и воспитания  относятся к разряду вечных, и пока живо человеческое  общество, они будут возникать, и станут предлагаться различные способы их разрешения. 

Естественно, что каждая эпоха заставляет ставить свои акценты на вопросах воспитания и обучения, постоянно пересматривать содержание обучения, приспосабливая его к нуждам общества, приводя его в соответствие с  научными взглядами и требованиями практики. Никому не придет в голову учить математике сейчас так, как это было, допустим, в конце XVIII века. Следует помнить, что жизнь быстро меняется, и за какие-нибудь тридцать — сорок лет во все области науки и прикладной деятельности вторглись вычислительные машины, в том числе малые. Скоро каждая домашняя хозяйка станет брать с собой в магазин  микрокалькулятор, а логарифмическая линейка станет  синонимом далекого прошлого. Можно ли в математике  продолжать обучение по-старому и не замечать этого нового,  властно вторгающегося в нашу повседневную жизнь? Очевидно, что нам придется пересмотреть наши педагогические  взгляды и на обучение таблице умножения, и на приобретение навыков пользования таблицами логарифмов или других функций. 

Процесс перестройки образования, особенно школьного, затрагивает интересы миллионов и очень сложен. К нему следует относиться с величайшей осторожностью, помня, что он сказывается на молодом поколении в течение всей жизни. Раны, наносимые психике учащихся при  непродуманном обучении, особенно тяжелы. Вот почему так  важно, чтобы реформаторы образования глубоко обдумывали не только формальную сторону содержания образования, но и его последствия. При этом ни на минуту нельзя забывать и того, что образование необходимо всем, а крупными специалистами в каждой данной области деятельности будут лишь единицы. Математические знания понадобятся всем, но математиками станут лишь немногие. Мы должны в  первую очередь помнить о нуждах всех, а не только тех, кто впоследствии станет участвовать в прогрессе  математической науки. 

С еще большей осторожностью следует относиться к  написанию учебников. Ведь любая ошибка —  математическая, педагогическая, психологическая — должна быть помножена на миллионы учащихся. Каждая неудачная формулировка обернется недопониманием сути дела  многими миллионами школьников и может оттолкнуть интерес их от изучаемого предмета. Следует, наконец, понять ту простую истину, что создание хорошего школьного  учебника требует не только огромного искусства, но и  предварительного общественного обсуждения, чтобы в нем не  проскочили отдельные ошибки, недомолвки, усложненные  определения и прочие недостатки. В создании учебника должны участвовать не только авторы и рецензенты, но и  многотысячная педагогическая общественность. 

Все новые педагогические идеи в жизнь школы вводит учитель, и поэтому он должен не только знать содержание и направление перестройки школьного образования, но и должен быть убежден в ее необходимости. Перед началом учебного года его нужно не только познакомить с  учебником, он должен убедиться, что учебник добротен, написан педагогически грамотно, учитывает психологические  возможности учащихся и соответствует педагогическим и  общественным установкам самого учителя. Конечно, было бы хорошо иметь не один, а два или даже несколько учебников по одному и тому же предмету, созданных на базе одной и той же программы, но исходящих из различных  педагогических установок. Тогда учителю будет легче найти  учебник, соответствующий его личным идеалам. Но и в этом случае учебники, прежде чем выйти миллионными  тиражами, должны пройти испытание педагогической  общественности. 

Настоящий выпуск составлен из работ известных  русских математиков-педагогов Н. И. Лобачевского, М. В.  Остроградского, А. Я. Хинчина. Все они были не только  вузовскими профессорами, но долгие годы занимались школьным просвещением. Н. И. Лобачевский более двадцати лет  занимал пост помощника попечителя огромного Казанского учебного округа; М. В. Остроградский в течение  десятилетий возглавлял математическое образование в русских  военных школах; А. Я. Хинчин на протяжении полутора  десятков лет возглавлял работы по математическому  образованию в Кабинете математики научно-исследовательского института школ Народного комиссариата просвещения РСФСР, а затем активно работал в Академии  педагогических наук РСФСР. 

В сборнике приводятся работы Н. И. Лобачевского, М. В. Остроградского и А. Я. Хинчина, посвященные  вопросам воспитания. Интересно отметить, что все эти работы роднит одна мысль, одно стремление — содействовать  развитию личности учащихся, способности самостоятельно мыслить, воспитанию высоких гражданских качеств, в том числе таких, как трудолюбие и стремление к познанию. Нам важно подчеркнуть, что каждый из названных трех авторов тесно увязывает проблемы образования с  проблемами воспитания. Их работы как бы перекликаются с  недавним постановлением ЦК КПСС «О дальнейшем  улучшении идеологической, политико-воспитательной работы»(«Известия», 7.V. 1979 г.). В этих работах русских математиков отстаивается та же самая мысль, которая столь четко  сформулирована в постановлении: «...добиваться органического единства учебного и воспитательного процессов,  формирования у учащихся и студентов научного мировоззрения, высоких морально-политических качеств, трудолюбия.  Прививать учащейся молодежи интерес к политическим  знаниям, всемерно развивать их общественную активность. Принять меры для дальнейшего развития внешкольной работы, технического и художественного творчества...». 

Само собой разумеется, что каждый автор, как бы ни был он велик, ограничен в своих суждениях  возможностями своего времени, а также теми проблемами, которые  выдвигала или выдвигает перед обществом реальная жизнь.

Ни Лобачевский, ни Остроградский не ставили перед собой ряда вопросов, которые теперь особенно актуальны и на которые обращено только что упомянутое  постановление ЦК КПСС. Они и не могли их ставить, поскольку жили в иную эпоху и их волновали другие проблемы. Но тем не менее многие освещенные ими вопросы не потеряли своей актуальности и в наши дни: чему и как учить, как  пробуждать интерес учащихся к обучению, как воспитывать  идеалы гражданственности, стремление к познанию и  творчеству, научное мировоззрение. Взгляды наших великих предшественников на эти вопросы нам интересны и сегодня. Вот почему мы объединили в одной брошюре три работы математиков-педагогов последних полутора столетий. 

Николай Иванович Лобачевский (1792—1856) был не только выдающимся математиком, но и замечательным  педагогом. Он начал свою педагогическую деятельность  двадцатилетним юношей и не прерывал ее до последних лет своей жизни. Его педагогические интересы не  ограничивались только чтением лекций, по своему положению  помощника попечителя огромного Казанского учебного округа он был близок к вопросам школьного образования и понимал, что образование и воспитание связаны между собой  неразрывными нитями. Недаром в первый год своего избрания на пост ректора Казанского университета Лобачевский  выбрал в качестве темы речи на торжественном собрании по случаю окончания учебного года и очередного выпуска окончивших обучение студентов вопросы воспитания. 

Лобачевским были затронуты в его речи действительно важнейшие вопросы воспитания: а) стремление к познанию, б) самостоятельность мышления, в) трудолюбие, г) научное мировоззрение, д) высокие культурные потребности, е)  убеждение в том, что знания должны быть обращены на пользу общества, ж) призыв к преподавателям помочь учащимся выявить и закрепить имеющиеся у них способности, з)  воспитание любви к отечеству и истинное понятие о чести. Для Лобачевского одного приобретения знаний еще  недостаточно, нужно воспитание всесторонне развитой личности, чтобы выпускник университета своими познаниями и  умениями составлял «честь и славу своего отечества». 

Речь «О важнейших предметах воспитания» впервые была опубликована в 1832 году в журнале «Казанский  вестник», ч. XXXV, кн. 8, с. 577—596. Сейчас она доступна читателям, поскольку в 1976 г. полностью была  опубликована в интересном томе: Н. И. Лобачевский.  Научно-педагогическое наследие. Руководство Казанским  университетом. Фрагменты. Письма. М., «Наука», 1976. 

В этом же томе содержатся , полезные материалы,  которые широко характеризуют Лобачевского как педагога и как крупного руководителя начального и среднего  образования. Он был активным сторонником развития поистине народного просвещения, а также женского образования. Будучи стихийным материалистом, Лобачевский стремился к тому, чтобы учащиеся получали широкое естественнонаучное образование и понимание основных законов  природы. 

В письме от 1.3. 1847 г. М. Н. Мусину-Пушкину —  попечителю Казанского учебного округа — Лобачевский  отмечает необходимость такого преподавания, которое  было бы понятно «для возраста», занимательно для учащихся, осуществляло бы «взаимную помощь и согласие разных  предметов», добивалось бы постепенности в обучении.  Одновременно в других документах он подчеркивал необходимость добиваться в первую очередь понимания учащимися  учебного материала и развития способности применять  полученные знания к решению частных задач. 

Чтобы составить себе хотя бы общее представление о размерах Казанского учебного округа, достаточно указать, что в него входили такие города: Казань, Нижний  Новгород, Пенза, Саратов, Астрахань, Пермь, Оренбург, Вятка. 

В письме от 22.4. 1846 г. директору училищ Саратовской губернии содержатся следующие интересные строки:  «усмотренная в некоторых учениках наклонность к  риторическим украшениям, нестрогости выражений и дикая  фантазия ... возлагают также на учителя словесности  обязанность заботиться, чтобы сочинения были писаны ясно,  изобиловали бы количеством мыслей, а не украшений, которые допускать только тогда, когда ими выражается особенная мысль или действительно усиливается выражение. Темы для сочинений, по преимуществу, должны касаться  истории, которой нравственное влияние: для молодых людей так полезно по примерам высоких добродетелей и  подвигов, увенчанных полезным успехом и награжденных  благодарностью современников и потомства». 

Лобачевский считал, что преподаватель имеет огромные возможности для привлечения внимания учащихся к  предмету обучения и что от его мастерства зависит многое в  достижении учебных успехов. В предисловии к своему  учебнику алгебры он писал: «Готов я думать, что если учение математики, столь свойственное уму человеческому,  остается для многих безуспешно, то это по справедливости  следует приписать недостаткам в искусстве и способе  преподавания» (Собр. соч., т. IV, с. 369). 

Михаил Васильевич Остроградский (1801—1862) —  выдающийся математик первой половины XIX столетия и  одновременно большой и разносторонний педагог, оказавший решительное влияние на подъем общего научного уровня в нашей стране. Его педагогические взгляды нашли исчерпывающее выражение в брошюре «Размышление о  преподавании», написанной им совместно с французским  педагогом И. А. Блумом и изданной в Париже в 1860 г. 

Эпоха шестидесятых годов вошла в историю России как время прогрессивных реформ и подъема русской культуры. Крымская война убедительно показала отсталость России и необходимость срочного развития просвещения народа и реорганизации всего образования. Все возникшие при этом вопросы и нашли отражение в книге «Размышления о преподавании». Ее авторы жестоко критиковали систему обучения, которое ведется слишком абстрактно и сухо, в отрыве от потребностей общества и интересов детей. 

Книга Остроградского и Блума была опубликована как раз в 1860 году и отражала насущные потребности  школьного образования как России, так и Франции. Оба автора едины в своем мнении, что преподавание оторвано от  потребностей жизни общества, оно слишком абстрактно и сухо и не учитывает наклонности и способности учащихся. Эту мысль они выразили значительно красочнее:  «Преподаватели гимназий, лицеев и военных школ признаются, что они читают лекции больше для скамеек и стульев, чем для внимательных и разумных учеников. Действительно, на уроках по арифметике, алгебре и геометрии ничто не  напоминает о насущной необходимости изучения этих предметов для практической жизни. Ничто не указывает на  наслаждение, испытываемое при изучении этих дисциплин  людьми, для которых это изучение связано с выбранной ими профессией. Ничего не рассказывают об истории наук. 

... глубокие сухие теории и непонятные определения формулируются, повторяются и пережевываются, давая только тот результат, что их усваивают очень небольшое число учеников». 

Авторы придают особое значение трудовому воспитанию детей с самого раннего возраста. По их мнению, дети,  начиная с семи лет, должны значительную часть времени  проводить в мастерских, которые необходимо создавать при каждой школе. Инструменты при этом следует поручать делать самим детям. Конечно, нет нужды переносить  буквально мысли и предложения Остроградского и Блума из XIX в. в наше время. Но сама идея воспитания  трудолюбия, ответственности к порученному делу, уважения к  физическому труду заслуживает внимания. 

Поразительно, что Остроградский, проживший всю жизнь в крепостнической России, нашел возможным сказать в адрес человека физического труда следующие полные  уважения слова: «... мы знаем огромное число людей без  литературного образования, которые не пишут ни поэм, ни  водевилей, но людей здравого смысла, которые возделывают поля, куют железо, перевозят на далекие расстояния  продукты земли, ткут волокно и шерсть, делают бумагу. 

Эти полезные люди могут сказать писателям: не очень-то удивительно, что вы знаете немного больше нас, вы, которых окружили заботами и обучали! Если бы нам  оказали десятую долю той помощи, которую вам расточали, мы были бы более полезными для себя и для других, но не были бы такими отчаявшимися и высокомерными, как вы!» 

Остроградский выступал как резкий противник раннего профессионализма, который не интересуется «ни вкусами, ни наклонностями» ребенка. Он возражает против желания многих готовить только выдающихся представителей науки, а также других сфер деятельности и призывает не  забывать о подготовке и воспитании «полезных и скромных  деятелей для наших современных обществ». 

Александр Яковлевич Хинчин (1894—1959) был не только выдающимся математиком, одним из создателей  теории случайных процессов и современной теории  вероятностей, но и первоклассным педагогом. Его лекции  воспринимались, с одной стороны, как глубокие философские  произведения, а х другой — как замечательные произведения искусства. Он никогда не стремился втиснуть в рамки одной лекции как можно больше фактического материала, а  старался всесторонне осветить основные идеи и понятия,  раскрыть перед слушателями душу предмета, его  содержание. На это он не жалел ни сил, ни времени. К студентам он относился с огромным уважением, видя в них своих будущих сотоварищей по работе. Он никогда не забывал о вопросах, которые ему задавали после лекций и отвечал на них либо немедленно, либо просил разрешения  предварительно обдумать ответ и в одну из ближайших встреч  приносил исчерпывающее разъяснение. К сожалению, имеется часть преподавателей, позволяющих себе обещать и  забывать свои обещания. Нет нужды говорить о том, что такое поведение оказывает резко отрицательное воспитывающее воздействие на учащихся. 

Хинчин считал, что преподаватель ответствен за все и должен быть исключительно строгим прежде всего к себе. Он не имеет права прийти на занятия недостаточно  подготовленным, не имеет права давать своим ученикам полузнани я, а должен в первую очередь во всех подробностях  выяснять принципиальные стороны предмета, его реальную и  логическую структуру. Те, кто слушал его лекции по курсу математического анализа, вспоминали о них и через много лет как о самых значительных университетских событиях, поскольку они не только представляли исключительное  математическое и художественное произведение, но и  приучали слушателей к мышлению, проникновению в само  существо предмета. Он любил читать лекции по  математическому анализу, поскольку именно этот предмет считал  основным для полноценного математического образования. Эти лекции А. Я. Хинчин читал для студентов механико-математического факультета Московского государственного университета и для студентов физико-математического  факультета педагогического института имени В. И. Ленина. Однажды он прочитал цикл лекций по математическому анализу для инженеров. Его курс собрал большую  аудиторию, которая не редела до конца цикла лекций. Основная трудность состояла в том, чтобы за короткий срок создать у слушателей концепцию предмета, внушить необходимость строгих доказательств для полноценных суждений и не допустимость поверхностных выводов. Этой цели он  добился с огромным успехом. 

При воспитании аспирантов Хинчин считал  необходимым направлять научные интересы своих учеников в самые актуальные направления научных поисков. В научной работе его привлекала не сложность использованного  аналитического и логического аппарата, а глубина результата, возможность получения из него разнообразных выводов как для самой математики, так и для различных ветвей  естествознания. Сам он в течение всей своей научной  деятельности живо интересовался статистической физикой и  проблемами телефонии. Его работы по телефонии привели к  созданию советской школы теории массового обслуживания. 

Приводимая в настоящем сборнике статья А. Я. Хинчина не была напечатана при жизни автора. Я обнаружил ее в его бумагах и подготовил к печати. Сначала она была опубликована в сборнике «Математическое просвещение» (М., Физматгиз, 1961, № 6, с. 7—28), а затем в 1962 г. в журнале «Математика в школе» (№ 3, с. 30—44). В том же 1963 г. эта статья была опубликована в специальном томе, подготовленном мной к изданию: А. Я. Хинчин.  Педагогические статьи. Изд-во Акад. пед. наук РСФСР, с. 128—160. 
 
На мой взгляд, статья Хинчина представляет  значительный интерес и теперь. Многие высказанные им положения требуют дополнительной разработки и существенных  дополнений, но тем не менее то, что уже сделано, по-прежнему интересно и полезно. Достаточно только перечислить те вопросы, которых коснулся автор: культура мысли  (правильность мышления, стиль мышления), моральные  моменты и воспитание патриотизма(честность и  правдивость, настойчивость и мужество, воспитание патриотизма). Одно это  перечисление показывает, что Хинчин в своей статье затронул наименее разработанные вопросы воспитания гражданина, которые можно и следует осуществлять на занятиях по  математике. 

Я убежден, что статья А. Я. Хинчина будет с пользой и интересом прочитана каждым читателем нашей серии. 

Проблемы воспитания и образования идут рядом и не могут быть разорваны. Преподаватель воспитывает не только тогда, когда он находится в классе и общается с  учениками, но всегда, где бы он ни находился. В декабре 1979 г. мне пришлось участвовать в сессии Ивановского университета, посвященной памяти крупного алгебраиста и увлеченного педагога академика Анатолия Ивановича Мальцева (1909—1967). Его ученики рассказывали о нем не только как о выдающемся математике и педагоге,  приучавшем к творческому восприятию математики, но и как о человеке, который всем своим поведением воспитывал студентов. В частности, сильное впечатление на  студентов производило то, что он, входя в университет, не  забывал поздороваться со швейцаром и с дежурной на вешалке. Этим он невольно подчеркивал уважение к любому виду труда, воспитывал в студентах мысль о человеческом  достоинстве каждого человека. Этим самым он разрушал  элементы гнусного сознания собственного превосходства не по делам, а по достигнутому общественному положению. 

Ни одна социальная система не может оставлять в  стороне, как нечто безразличное, вопросы воспитания  подрастающего поколения в необходимом обществу духе. 

Работая в Московском государственном университете и встречаясь с молодыми людьми из стран, только что  освободившихся от колониального рабства, мне  неоднократно приходилось сталкиваться с одной особенностью  характера их обучения. Молодые люди, закончившие там  университеты, приучались к тому, чтобы становиться послушными исполнителями, а не творцами. Мне приходилось затрачивать огромные усилия на то, чтобы пробудить в них  стремление не просто заучить познанное, а критически его  осмыслить, выяснить возможности дальнейшего продвижения на пути познания пока еще неизвестного. Года два проходило, прежде чем молодые люди приучались самостоятельно оценивать прочитанное и смотреть на него собственными глазами. 

Перед Великой Октябрьской революцией церковь и  царское правительство внушали русским их превосходство  перед «инородцами». К счастью, мы освобождены теперь от этого страшного наследия прошлого и прекрасно знаем, что любой народ способен дать и дает таланты в любых видах деятельности: науке, строительстве, промышленности,  экономике, военном деле, поэзии и пр. Не принадлежность к той или иной социальной прослойке, к той или иной расе, а условия и труд создают народные способности. Рабство, колониализм, социальная несправедливость — вот  источники замедленного умственного развития. Каждый народ, каждая социальная группа, когда они находятся в  благоприятных условиях, способны на высшие проявления  духовной деятельности. 

Просвещению народа, развитию его культуры и науки, воспитанию трудолюбия, ответственного отношения к  труду, духа интернационализма и одновременно советского патриотизма, научного мировоззрения Советское  государство уделяло исключительное внимание с первых дней своего существования. Впервые в истории человечества для всех граждан, независимо от социального положения и  национальной принадлежности, открылись двери школ и  университетов, образование стало бесплатным. Результаты не замедлили сказаться — за короткий исторический срок в нашей стране были воспитаны миллионы  квалифицированных рабочих, тысячи талантливых ученых, инженеров, изобретателей, педагогов, врачей, деятелей культуры. Возникла мощная промышленность, созданная разумом и руками народа. Все народы и народности страны  оказались вовлеченными в этот бурный процесс созидания, внесли свою лепту в развитие науки, культуры и экономики государства. 

Само собой разумеется, что с развитием общества как проблемы образования, так и проблемы воспитания  подрастающего поколения- не остаются неизменными: одни проблемы отмирают или становятся не столь существенными, другие же возникают вновь или же им приходится придавать большее, чем раньше, значение. 

Именно по этой причине время от времени партия и  правительство обращают внимание на актуальные вопросы воспитания. За последние годы были изданы: 

постановление ЦК КПСС и Совета Министров СССР «О дальнейшем совершенствовании обучения, воспитания учащихся общеобразовательных школ и подготовки их к ТРУДУ» (1977 г.); 

постановление ЦК КПСС «О дальнейшем улучшении идеологической, политико-воспитательной работы» (1979 г.); 

постановление ЦК КПСС и Совета Министров СССР «О дальнейшем развитии высшей школы и повышении  качества подготовки специалистов» (1979 г.). 

В каждом из этих постановлений вопросы воспитания занимают значительное место и среди них особую роль играют вопросы воспитания научного мировоззрения на всех этапах развития до школы, в школе, в ВУЗе и на  работе. 

Проблема воспитания научного мировоззрения  исключительно широка и разнообразна. Она включает в себя не только создание общественных и политических идеалов, но и воспитание ответственного отношения к труду,  порученному делу, к качеству выполнения работы, развития честности, благородства, стремления к познанию нового, духа творческих исканий. Сюда же относится воспитание травильных взглядов на окружающий нас мир, присущие ему закономерности, происхождение научных знаний и научных понятий, назначение науки, взаимоотношение теории и практики. Далее к сказанному следует добавить необходимость воспитания веры в собственные силы, в  способности самостоятельно преодолевать трудности и творить. Но эта вера должна базироваться не на неумеренном  самомнении, а на реальном умении познавать и мыслить, умении напряженно работать и добиваться поставленных перед собой целей. 

Воспитание научного мировоззрения у подрастающего поколения представляет собой одну из важнейших задач, стоящих перед советской школой. Оно требует постоянной, дружной и согласованной работы всего педагогического коллектива школы, преподавателей всех дисциплин.  Математика не составляет здесь исключения. Воспитание  научного мировоззрения должно быть не навязчивым и учащиеся как бы сами должны подходить к восприятию  определенных взглядов, норм поведения, взаимоотношению с коллективом, отношению к труду. Всякая попытка  административным путем заставить учащихся придерживаться определенных взглядов, если только они не стали их  внутренним убеждением, обречена на неудачу. Здесь особенно важно подчеркнуть, что личный пример учителя при этом играет решающую роль. Подростки все воспринимают  прямолинейно и любая неискренность учителя оставляет в их сознании долгие годы незаживающую рану. 

Стремление к совершенству, к отличному выполнению порученного дела приходит не сразу, а должно  воспитываться с раннего детства и становиться частью  мировоззрения. Нетерпимость к небрежному выполнению  обязанностей, как бы малы они ни казались, должна войти в плоть и кровь. Очень важно создать привычку делать каждое дело хорошо сразу, без черновика. Дело в том, что в наши дни появилось большое число профессий, которые требуют мгновенных решений и тщательного продумывания всех «мелочей» еще до начала выполнения задания. Всякая ошибка может привести к катастрофическим последствиям. Хирург, оперирующий на сердце, не может допустить ни одного ошибочного движения: его ошибка ставит пациента на грань гибели. Летчик в сложной ситуации обязан  принимать немедленное решение; он не имеет права и времени сделать необходимое начерно. Ошибка, допущенная  диспетчером аэропорта или железной дороги, может быть  причиной непоправимой катастрофы. Вот почему так важно к концу школьного обучения не только не приучать к  черновикам, а отучать от них, приучать к дисциплине  мышления и отучать от мысли, что недоделки и ошибки удастся устранить в чистовике. 

Другое дело стремление к совершенству, повышенной требовательности в первую очередь к самому себе, поиску лучшего, более приспособленного к требованиям жизни. Это весьма важный момент в воспитании мировоззрения. Наша задача состоит в том, чтобы последующее поколение стало поколением творцов, ищущих лучшие подходы к  решению возникающих проблем, находящих более  совершенные формы труда, увлеченных выполняемым делом,  получающих радость от превосходно выполненного задания, стремящихся приносить добро людям и заботящихся об окружающей нас природе. 

Математические знания в руках мыслящих и ищущих людей способны творить чудеса. Если говорить о  производстве, то только математика позволяет найти оптимальные способы расположения цехов на заводе, станков внутри цеха, когда достигаются минимальные расходы на  транспортировку полуфабрикатов. Сейчас исключительно остро стоит задача управления качеством продукции в процессе производства. Она приводит к интересным и сложным математическим вопросам. 

Очень важно воспитывать каждого школьника в  убеждении важности математики и ее методов для жизни  современного общества. Воспитать это убеждение одними словами невозможно, необходим показ математики в действии. Только так можно добиться, чтобы школьник видел за вычислениями и формальными преобразованиями, за  геометрическими образами не только абстрактные символы и «игру ума», но и реальное содержание и метод познания, пусть даже в самой простой форме. 

Учащиеся должны постоянно видеть, что  математические понятия, методы и теории возникли не в результате простого желания ученых, — они были подготовлены  насущными практическими потребностями, и нужно было длительное время, чтобы первичные смутные  представления превратились в понятия и твердо установленные  математические факты. Поскольку жизнь общества  развивается, а не стоит на месте, постоянно возникают новые - вопросы, для которых математика должна разрабатывать также новые приемы решения и такие понятия, в терминах которых удается, как говорят математики, поставить  задачу: сформулировать ее точно. Но для этого абсолютно  необходимо, чтобы учитель был вооружен, с одной стороны, знанием истории собственной науки, а с другой — имел в руках материал, который позволял бы ему в доступной форме излагать сведения о развитии математики в наши дни, о ее многочисленных и неразрывных связях с  практикой во всем ее многообразии. 

То, о чем мы сейчас говорим, является не чем иным, как частью воспитания научного мировоззрения. Этим самым я настаиваю на том, что воспитание научного мировоззрения должно начинаться с раннего детства и продолжаться всю жизнь. Пусть в школе для этого имеются ограниченные возможности (но все же немалые), учитель математики не может отказываться от формирования научного  мировоззрения на своих уроках. 

Научное мировоззрение формируется при усвоении представления о единстве мира, его существовании вне нашего сознания и познаваемости его законов. Математика в этом отношении может дать многое. То обстоятельство, что одни и те же математические результаты применимы не только к разнообразным явлениям природы, а и к  производственным, экономическим и техническим процессам,  является веским свидетельством общности количественных закономерностей, свойственных столь различным областям знания. Необходимость усвоения идеи универсальности законов природы не уставал подчеркивать в своих трудах В. И. Ленин. И в процессе установления этого единства математике он отдавал должное. 

Естественно, что никак нельзя терять те возможности, которые имеет школа. А возможности эти совсем не малы — они включают и межпредметные связи, и выбор системы  материала в учебнике, и умелое привлечение  иллюстративных примеров. 

Математика обладает крайней абстрактностью своих понятий и одновременно исключительной широтой  применимости результатов к разнообразным областям знания и деятельности. В чем здесь причина? В чем состоит сила абстракции? Почему абстрактная теория позволяет  улавливать количественные закономерности реальных явлений, притом столь различных по своей физической природе? Как происходит формирование понятий математики и  возникают ее теории? Откуда берутся новые направления  математических исследований? Все эти вопросы требуют  ответа и притом ответа своевременного и доступного  пониманию обучающихся. Удачное рассмотрение их педагогом, а также отсылка к хорошо написанным на эти темы  брошюрам и книгам позволит внести серьезный вклад в  воспитание диалектического мировоззрения. 

Важность для повседневной жизни и практики таких привычных понятий, как число и прямая, не вызывают ни у кого сомнений. История их становления способна не только вызвать живой интерес учащихся, но и помочь  формированию диалектического мировоззрения. Первые шаги в развитии счета, проиллюстрированные отрывками из произведений далекого прошлого, рассказов этнографов и путешественников, дают нечто большее, чем только изложение этапов пути развития числа, способное  заинтересовать учащихся. Они создают основу исторического подхода к пониманию процесса возникновения и развития научных понятий. Кроме того, на этом примере учащиеся наглядно убеждаются в том, что практика во всем ее  разнообразии является основной движущей силой прогресса  математики (как, впрочем, и остальных наук). Конечно,  понятие числа служит лишь одним из примеров, которые  способен рассмотреть учитель в классе. 

Математическая символика, играющая столь большую роль в науке, позволяет не только компактно записывать информацию о задаче, но и автоматизировать ее решение. Достаточно вспомнить решение систем линейных  алгебраических уравнений, чтобы убедиться в правильности этого утверждения. 

История науки играет исключительно важную роль в формировании мировоззрения учащихся в самых  разнообразных планах. Мне дополнительно к тому, что уже было сказано, хотелось бы отметить ее значение для воспитания патриотизма и интернационализма. Наш народ внес  заметный вклад в прогресс математики и ее применений к  вопросам естествознания, инженерного дела и других  областей знания. У нас есть чем и кем гордиться. Имена М. В. Остроградского, Н. И. Лобачевского, П. Л. Чебышева, А. М. Ляпунова, А. А. Маркова, А. Н. Крылова, Н. Е. Жуковского, К. Э. Циолковского, С. Н. Бернштейна, А. Я. Хинчина и многих других прочно вошли в историю мировой науки. С этими именами связаны многочисленные фундаментальные достижения математики — решение  ряда задач, стоявших перед наукой длительное время,  создание новых направлений исследований, коренной сдвиг применений математики к решению принципиальных  вопросов практики. 

М. В. Остроградский внес большой вклад в решение фундаментальных задач механики и математической  физики. Ему же принадлежат первые попытки исследования  вопроса о контроле качества массовой промышленной  продукции с помощью методов теории вероятностей. Глубокие идеи Н. И. Лобачевского о геометрической структуре  пространства и о логической возможности иной, чем  евклидова, геометрической системы, а также мысли о том, что явления макро- и микромира могут требовать для своего описания различных геометрий, опередили свою эпоху. 

Работы П. Л. Чебышева внесли фундаментальный вклад в развитие теории чисел, теории вероятностей, теории  механизмов и привели к формированию новой ветви  математики — теории аппроксимации функций. 
 
А. М. Ляпунов внес исключительный по значимости вклад в создание теории устойчивости движения, в изучение фигур равновесия вращающейся жидкости, в теорию  вероятностей. Его результаты и идеи не потеряли своего  значения и до настоящего времени и продолжают интенсивно разрабатываться как в теоретическом, так и прикладном плане. 

А. А. Марков создал себе имя работами по теории  квадратичных форм и теории вероятностей. Начатая им область исследований в теории вероятностей получила название цепей Маркова. В свою очередь это дало толчок к  позднейшему созданию теории случайных процессов,  являющейся теоретической основой ряда физических явлений и многих процессов живой природы, инженерного дела, экономики. 

Многие учащиеся прекрасно знают, что с именем Н. Е. Жуковского связана целая эпоха в развитии науки — создание теории полета самолета, методы расчета  подъемной силы крыла, решение многочисленных задач гидро- и аэродинамики. Недаром В. И. Ленин назвал Жуковского отцом русской авиации. 

Подавляющее большинство школьников прекрасно знает, что с идеями К. Э. Циолковского связано создание новой ветви науки — математической космонавтики и что ему принадлежит ряд конструктивных идей, которые  широко используются в настоящее время как при  конструировании ракет, так и при выборе методов их расчета. 

То, что я сказал, хорошо известно многим учителям  математики, но речь идет о том, что, сообщая школьникам необходимые факты истории в нужное время, можно тем самым добиться решения по меньшей мере двух задач: воспитания патриотизма, гордости за величайшие  достижения мысли нашего народа, с одной стороны, и с другой — воспитания уверенности учащихся в том, что они сами  обладают творческими способностями, которые следует  выявись в напряженном труде. 

Конечно, было бы глубокой ошибкой ограничиваться беседами только об отечественных ученых и их вкладе в прогресс науки. В развитие науки вносят свой вклад все народы мира, и нет какого-то избранного народа, который обладает особой талантливостью. Мы знаем теперь, как велик вклад в прогресс математических и  астрономических знаний народов Средней Азии. Разрушительные войны и нашествие орд завоевателей уничтожили культурные традиции и создали длительный застой в научной  жизни этих народов. Советская власть пробудила творческие силы узбеков, таджиков, казахов, киргизов, туркмен, и мы можем наблюдать, сколь велики успехи народов Средней Азии во всех проявлениях интеллектуальной деятельности. 

Я убежден, что для помощи педагогу в организации  воспитательной работы, в том числе и в воспитании научного мировоззрения, необходимо создать соответствующие  пособия. Естественно, нужно, чтобы эту работу возглавили учреждения Академии педагогических наук СССР. Эти  книги должны быть такими, чтобы учитель мог легко  использовать прочитанное для своей повседневной работы. 

Мировоззрение человека возникает не вдруг, а  постепенно, в результате размышлений и анализа поступков  собственных и других лиц, в результате размышлений над  услышанным, прочитанным и наблюденным. Задача учителя помочь учащимся встать на правильный путь. Хорошо выполненная работа должна систематически отмечаться, так же как плохо выполненная работа, нерадивое отношение к порученному делу должно находить соответствующую оценку. Отклонение от такой позиции исключительно  опасно в социальном отношении, поскольку оно воспитывает потребительское отношение, наводит на мысль о том, что кто-то обязан работать на бездельника, а он может  пользоваться плодами труда других. 

Совершенствование обучения и воспитания молодого поколения крайне необходимо всему народу и каждой  отдельной семье. Это сложнейший процесс, в котором  участвуют многие сотни тысяч педагогов. Удачные решения этих вопросов должны быть сделаны достоянием всех  учителей. 

Особенно важно воспитывать в каждом ученике  нетерпимость к формализму знаний, когда заученные сведения не подкрепляются умением активно их применять при  необходимости. Формализм математических знаний является тяжелой и, к сожалению, еще весьма распространенной болезнью. А. Я. Хинчин дал в одной из своих статей{\footnote{Хинчин А. Я. О формализме в школьном преподавании математики. — «Советская педагогика», 1944, № 11—12, с. 21—27.}} четкое определение этого явления: «Для всех проявлений формализма характерно неправомерное доминирование в сознании и памяти учащихся привычного внешнего (словесного, символического или образного) выражений  математического факта над содержанием этого факта». 

Как часто мы сталкиваемся с этим явлением не только в школе, но и вузе, когда существо понятия или результата остается вне сознания учащегося и он подменяет истинное понимание его суррогатом — бездумным заучиванием.  Содержание математического понятия или математического факта не доходит до сознания учащегося, в памяти остаются лишь неосознанные слова. Естественно, что при таком  изучении математики нельзя говорить не только о применении ее положений к реальным задачам, но и о дальнейшем ее изучении. При этом не развивается мышление и  способность самостоятельного подхода к новым задачам.  Формализм математических знаний, допущенный на каком-либо одном этапе познания, тяжелым грузом сопровождает учащегося на всем продолжении дальнейшего обучения. Вот почему так важно Строить обучение таким образом, чтобы не давать возможности появиться формализму  знаний в зародыше. Необходимо, чтобы преподаватель  постоянно следил за этим, чтобы система разработанных им  вопросов и задач позволяла своевременно обнаруживать, а затем и способствовать уничтожению этих бацилл формализма. 

Отсутствие формализма знаний является непременным условием полноценного раскрытия способностей и  творческого начала. Развитие творческих способностей требует длительных усилий педагогического коллектива и  активного содействия самих учащихся. Воспитанию стремления к творчеству следует уделять пристальное внимание на всех этапах обучения. Каждый предмет школьного круга может оказать свое неповторимое воздействие на  формирование творческого облика учащегося. Математика  обладает для этого исключительными возможностями.  Действительно, поиск решений нестандартных задач,  размышления над парадоксами, анализ содержания условий теорем и сути их доказательств, беседы о путях работы  творческой мысли выдающихся ученых — все это составляет вехи на пути развития способностей и самого духа  творческого горения. 

Чем раньше учащийся почувствует, что он обладает  задатками замечательного дара — самостоятельностью мышления при поиске нового, тем лучше будет для него самого и для общественного прогресса. Но при этом  
обязательно следует выработать и привычку к  систематической работе, поскольку без регулярного напряженного труда никакие потенциальные способности не могут  привести к успеху. Хороший учитель сумеет тактично внушить своему ученику, что наличие творческого начала в какой-либо области знаний не дает права без должного уважения относиться к другим учащимся, которые не смогли проявить своих способностей в данной области деятельности. Они могут иметь способности и таланты совсем другого рода, которые также необходимы обществу. Неуважения  заслуживают только безделие и лень мысли. 

Всестороннее развитие учащихся, выявление их  способностей, показ математики во всем ее разнообразии — больше и крайне важные для общественного развития  задачи. В их решении нас ждет увлекательная, трудная и  исключительно важная работа, за которую следует  приниматься немедленно. Нет сомнений, что предлагаемая  брошюра внесет свою лепту в решение проблемы воспитания молодого поколения. 


	\section[\textls{Н. И. Лобачевский}. Фрагменты речи «О важнейших предметах воспитания»]{{\small Н. И. ЛОБАЧЕВСКИЙ} \\ \textbf{ФРАГМЕНТЫ РЕЧИ «О ВАЖНЕЙШИХ ПРЕДМЕТАХ ВОСПИТАНИЯ»}{\footnote{Лобачевский Н. И. Научно-педагогическое наследие. Руководство Казанским университетом. Фрагменты. Письма. М., «Наука», 1976, с. 16—21. }}}

«Воспитание ... начинается от колыбели, приобретается сперва одним подражанием, постепенно развертывается ум, память, воображение, вкус к изящному, пробуждается любовь к себе, к ближнему, любовь славы, чувство чести, желание наслаждаться жизнью. Все способности ума, все дарования, все страсти, все это обделывает воспитание, соглашает в одно стройное целое, и человек, как бы снова родившись, является творение в совершенстве. 

Наружный вид его, возвышенное чело, взор, который всюду устремляется, все созерцает вверху, вокруг себя; черты лица, в которых изображается чувственность,  покорная уму, — все показывает, что он родился быть  господином, повелителем, царем природы. Но мудрость, с  которой он должен править с наследственного своего престола, не дана ему от рождения: она приобретается учением. 

В чем же должна заключаться эта мудрость? Чему  должно нам учиться, чтобы достигнуть своего назначения? Какие способности должны быть раскрыты и  усовершенствованы, какие должны потерпеть перемены; что надобно придать, отсечь, как излишнее, вредное? 





Мое мнение: ничего не уничтожать и все  
усовершенствовать... Всего обыкновеннее слышать жалобы на страсти, 
но, как справедливо сказал Мабли 3: чем страсти сильнее, 
тем они полезнее в обществе; направление их может быть 
только вредно. 
Что же надобно сказать о дарованиях умственных,  
врожденных побуждениях, свойственных человеку желаниях? 
Все должно остаться при нем; иначе исказим его природу, 
будем ее насиловать и повредим его благополучию. 
Обратимся, во-первых, к главнейшей способности, уму, 
которым хотят отличить человека от прочих животных, 
противуполагая в последних инстинкт. Я не того мнения, 
чтобы человек был лишен инстинкта, который является 
во многих действиях ума, который в соединении с умом  
составляет Гений. Замечу только мимоходом, что инстинкт 
не приобретается; Гением быть нельзя, кто не родился. 
В этом-то искусство воспитателей: открыть Гений,  
обогатить познаниями и дать свободу следовать его внушениям...» 
«Как бы то ни было, но в том надобно признаться, что не 
столько уму нашему, сколько дару слова одолжены мы всем 
нашим превосходством пред прочими животными. Из них 
самые бликие по сложению тела своего, как уверяют анато- 
мики, лишены органов, помощию которых могли бы  
произносить сложные звуки. Им запрещено передавать друг 
другу понятия. Одному человеку предоставлено это право; 
он один на земле пользуется сим даром; ему одному велено 
учиться, изощрять свой ум, искать истин соединенными 
силами. Слова, как бы лучи ума его, передают и  
распространяют свет учения. Язык народа — свидетельство его 
образованности, верное доказательство степени его  
просвещения. 
Чему, спрашиваю я, одолжены своими  
блистательными успехами в последнее время математические и физиче- 
3 Габриель Бонно де Мабли (1709—1785) —  
представитель французского просвещения, историк и политический  
мыслитель, развивал идеи утопического коммунизма. Лобачевский, по- 
видимому, имеет в виду его книгу по этике «Principes de morale».— 
24 
ские науки, слава нынешних веков, торжество ума  
человеческого? Без сомнения, искусственному языку своему, ибо 
как назвать все сии знаки различных исчислений, как не 
особенным, весьма сжатым языком, который, не утомляя 
напрасно нашего внимания, одной чертой выражает  
обширные понятия. Такие успехи математических наук,  
затмивши всякое другое учение, справедливо удивляют нас; 
заставляют признаться, что уму человеческому  
предоставлено исключительно познавать сего рода истины, что он, 
может быть, напрасно гоняется за другими; надобно  
согласиться и с тем, что математики открыли прямые средства к 
преобретению познаний. Еще не с давнего времени  
пользуемся мы сими средствами. Их указал нам знаменитый 
Бэкон 4 . Оставьте, говорил он, трудиться напрасно,  
стараясь извлечь из одного разума всю мудрость; спрашивайте 
природу, она хранит все истины и на вопросы ваши будет 
отвечать вам непременно и удовлетворительно. Наконец, 
Гений Декарта5 привел эту счастливую перемену, и,  
благодаря его дарованиям, мы живем уже в такие времена, 
когда едва тень древней схоластики бродит по  
университетам. Здесь, в это заведение вступивши, юношество не 
услышит пустых слов без всякой мысли, одних звуков без 
всякого значения. Здесь учат тому, что на самом деле  
существует, а не тому, что изобретено одним праздным умом. 
Здесь преподаются точные и естественные науки с пособием 
языков и познаний исторических. Здесь преподаватели  
разделяют между собою предметы, которыми всю жизнь свою 
занимаются, еще с молодых лет почувствовав в себе охоту 
и некоторые дарования. Как жалко, что истинному  
просвещению предпочитаются суетные выгоды домашнего  
воспитания. Кто хочет образовать своих детей для государства, 
тот должен прибегнуть к средствам, которые одно только 
государство в состоянии доставить; тот должен учить своих 
детей в общественных заведениях. 
Одно воспитание умственное не довершает еще  
воспитание. Человек, обогащая свой ум познаниями, еще должен 
учиться уметь наслаждаться жизнию. Я хочу говорить об 
образованности вкуса. 
4 Бэкон Фрэнсис (1561 — 1626). Новый органон. Л., 1935, 
с. 384. Ф. Бэкон стремился освободить науку от влияния теологии 
и схоластики.— Б. Г. 
6 Декарт Рене (1596—1650). Рассуждение о методе, с  
приложениями. М., Изд-во АН СССР, 1953, с. 656.— Б. Г. 
25 
Жить — значит чувствовать, наслаждаться жизнию, 
чувствовать непрестанно новое, которое напоминало бы, 
что мы живем. Так, стихотворец наш Державин говорит о 
ЛЮДЯХ! 
Непостоянство — доля смертных 
В пременах вкуса — счастье их. 
Среди утех своих несметных 
Желаем мы утех иныхб, 
о.Подобно реке она [жизнь] течет в излучистых  
берегах: то разливается в лугах радости, то омывает крутые 
утесы горестных размышлений. Ничто так не стесняет сего 
потока, как невежество: мертвою, прямою дорогою  
провожает оно жизнь от колыбели к могиле. Еще в низкой доле 
изнурительные труды необходимости, мешаясь с  
отдохновением, услаждают жизнь земледельца и ремесленника; 
по вы, которых существование несправедливый случай  
обратил в тяжелый налог другим; вы, которых ум отупел 
и чувство заглохло; вы не наслаждаетесь жизнию. Для вас 
мертва природа, чужды красоты поэзии, лишена прелести 
и великолепия архитектура, незанимательна история веков. 
Я утешаюсь мыслию, что из нашего университета не выйдут 
подобные произведения растительной природы; даже 
не войдут сюда, если, к несчастию, уже родились с 
таким назначением. Не выйдут, повторяю, потому, что 
здесь продолжается любовь славы, чувство чести и  
внутреннего достоинства. 
Кажется, природа, одарив столь щедро человека при его 
рождении, еще не удовольствовалась. Вдохнула в каждого 
желание превосходить других, быть известным, быть  
предметом удивления, прославиться и, таким образом,  
возложила на самого человека попечение о своем  
усовершенствовании...» 
«Будем же дорожить жизнию, покуда она не теряет 
своего достоинства. Пусть примеры в истории, истинное 
понятие о чести, любовь к отечеству, пробужденная в юных 
летах, дадут заранее то благородное направление страстям 
и ту силу, которая позволяет торжествовать над ужасом 
смерти». 
«...Вы счастливее меня, родившись позже...  
счастливейшие дни России еще впереди. Мы видели зарю,  
предвестницу их, на востоке; за нею показалось солнце... Я все 
сказал этим». 
6 Державин Г. Р. Из оды «К первому соседу». См.! 
Державин, Стихотворения. Л., 1933, с. 151 — 152.— Б. Л 
26 
Н. И. ЛОБАЧЕВСКИЙ 
НАСТАВЛЕНИЕ УЧИТЕЛЯМ МАТЕМАТИКИ 
В ГИМНАЗИЯХ, СОСТАВЛЕННОЕ 
г. РЕКТОРОМ УНИВЕРСИТЕТА 
ОРДИНАРНЫМ ПРОФЕССОРОМ 
ЛОБАЧЕВСКИМ 
16 АВГУСТА 1830 г.7 (ФРАГМЕНТЫ? 
«В математике всего важнее способ преподавания. 
Обширность науки даже в первых ее началах,  
которые должны составлять гимназическое учение, уже  
такова, что может быть обнимаема только в общих  
правилах. Чтобы прийти к сим правилам, надобно частные и 
раздельные представления о мере и числе соединить в 
одно и с такими-то сложными и отвлеченными понятиями 
рассуждать о всяком предмете и счете. Ясность предмета 
и порядок, в котором строгое суждение связывает все  
истины, служат единственным средством, чтобы постигнуть 
и удержать общие правила. Польза от сего учения бывает 
двоякая: применение его к потребностям нашей жизни 
и дальнейшее развитие самой науки. Умение применять 
общие правила ко всякому случаю предполагает твердое 
их познание и, сверх того, навык». 
«...в постепенном развитии понятий и в умении не  
допускать, чтобы одно изучение на память общих правил 
и механическое исчисление заменяли суждение,  
заключается искусство преподавания и успех его». 
«...математическим наукам служат те первые понятия, 
которые мы получаем в природе прямо чувствами; даже 
первые наши суждения о предметах, составляющих сии 
понятия, заключаются более в чувствах по навыку, нежели 
в действии ума, когда он под общим видом обнимает все 
возможные случаи». 
«...весьма важно условие для развития способностей 
то, чтобы учение не было механической работой и чтобы 
ученик постигал чувствами то, чего не в состоянии  
постигать суждением. Другое необходимое условие для  
успешного учения то, чтобы уметь победить леность и  
рассеянность детского возраста; той и другой цели достигает 
7 Лобачевский Н. И. Научно-педагогическое наследие, 
с. 626-528. 
27 
в совершенстве способ взаимного обучения, который  
разнообразием своим предохраняет детей от скуки». 
Н. И. ЛОБАЧЕВСКИЙ 
ИЗ ПРЕДПИСАНИЯ ДИРЕКТОРУ 
ПЕНЗЕНСКОГО ДВОРЯНСКОГО ИНСТИТУТА 
О СПОСОБАХ УЛУЧШЕНИЯ 
ПРЕПОДАВАНИЯ МАТЕМАТИКИ. 
5 МАРТА 1847 г.8 
...преподавание математики бывает только тогда  
успешным, когда ученики вполне понимают учителя, а потому 
должен он приспособиться к понятию слушателей,  
присоединять занимательность к своему преподаванию и не  
спешить вперед идти, покуда ученики не будут в состоянии 
за ним следовать. Занимательность учения заключается 
в удовольствии понимать предмет и преподанное  
применять к решению вопросов. Учитель должен за решением 
следить, руководствовать и каждого ученика в его хороших 
понятиях одобрять». 
М. В. ОСТРОГРАДСКИЙ и А. БЛУМ 
РАЗМЫШЛЕНИЯ О ПРЕПОДАВАНИИ 
(ФРАГМЕНТЫ)9 
«...для обучения молодежи используют те же приемы,  
которыми пользовались Сократ и Платон для преподавания 
высоких истин морали людям, уже сформировавшимся  
путем глубокого изучения логики и философии, людям,  
изощренным в искусстве речи». 
«В течение долгого времени эти абстрактные методы 
были единственным способом, которым пользовались для 
распространения научных знаний. 
8 Лобачевский Н. И. Научно-педагогическое наследие, 
с. 533. 
9 Остроградский М. В. Педагогическое наследие. 
Документы о жизни и деятельности. М., Гос. изд-во физ.-мат. лит., 
1961, с. 32-52. ' 
28 
Впрочем, надо отметить, что наука, в современном  
понимании, имеет за собой весьма небольшое число веков. 
Поэтому не удивительно, что так медленно развивались 
методы обучения». 
«Нам кажется, что после изобретения письменности 
самым большим открытием было использование  
человечеством так называемой десятичной системы счисления. 
Мы хотим сказать, что соглашение, с помощью которого 
мы можем выразить все полезные числа двенадцатью  
словами и их окончаниями, является одним из самых  
замечательных созданий человеческого гения...» 
«Заметим мимоходом, что никто не приводит такие  
исторические соображения, а они ведь так хорошо помогают 
привлечь внимание аудитории к урокам преподавателя. 
Но сразу ли дало результаты это замечательное  
открытие, сделанное при зарождении общества? 
Нет. На часах истории потребовалось пятьдесят веков, 
чтобы прийти к тому замечательному способу, которым мы 
обладаем для записи чисел. 
Всего около девяти веков прошло с тех пор, как мы 
научились записывать числа с помощью цифр, каждая 
из которых имеет свое собственное значение и значение, 
зависящее от ее положения. 
Этот принцип относительно цифр очень прост и, однако, 
только, можно сказать, случайно он стал общепринятым, 
притом очень медленно, в Европе и в остальном мире. 
Нам кажется, что важность такого глубокого открытия 
подчеркивают недостаточно и внимания ему уделяют  
слишком мало. 
Действительно, какие вычисления были возможны до 
этого открытия? Все тормозилось из-за отсутствия такой 
простой записи чисел. 
Все достижения математических наук, астрономии,  
механики, даже химии зависели от выполнения в.уме  
чрезвычайно сложных действий. 
Ныне десятилетний ребенок может без труда выполнить 
вычисления, которые даже не могли себе представить  
великие Архимед, Пифагор или Гиппарх. 
Таким образом, арифметика, скромная арифметика, 
является относительно недавним открытием. Теперь она 
удивительно проста, если ее не усложняют для забавы 
педантическими ухищрениями». 
«...прогресс прикладных наук, делая ощутимой  
необходимость в подготовке инженеров, офицеров, медиков, 
29 
тактиков, промышленников, требует изменения системы 
преподавания. Стало понятным, что техническое обучение 
детей следует начинать в более раннем возрасте». 
«...кто из нас не испытал несказанных огорчений в  
самом начале обучения элементам наук? 
Но мало было понять, следовало еще и запомнить то, 
что усваивалось с неслыханными трудностями. Кто из 
нас не видел, что из пятидесяти соучеников по меньшей 
мере сорок испытывали отвращение и падали духом из-за 
абстрактности идей, преподносимых нам до того, как они 
становились понятными на примерах, взятых из  
житейской практики? 
Преподаватели гимназий, лицеев и военных школ  
признаются, что они читают лекции больше для скамеек и 
стульев, чем для внимательных и разумных учеников. 
Действительно, на уроках по арифметике, геометрии 
и алгебре ничто не напоминает о насущной необходимости 
изучения этих предметов для практической жизни. Ничто 
не указывает на наслаждение, испытываемое при изучении 
этих дисциплин людьми, для которых это изучение связано 
с выбранной ими профессией. Ничего не рассказывают об 
истории наук. 
Осмелимся заявить, что глубокие сухие теории и  
непонятные определения формулируются, повторяются и  
пережевываются, давая только тот результат, что их  
усваивает очень небольшое число учеников. 
Создается впечатление, что тайнами науки все еще 
владеют служители культа Древнего Египта». 
«Кто сможет отрицать, что имеет первостепенное  
значение упрощение методов, чтобы столь желательное  
приобщение к наукам сделать легким для преподавателя и  
приятным для ребен ка?» 
«Мы без колебания заявляем, что изучение биографий 
людей, принесших пользу наукам и искусству, является 
одним из средств, которые мы используем, чтобы привлечь 
внимание учеников. Это в одно и то же время отличная 
разрядка и средство с помощью живого рассказа  
запечатлеть то или иное основное положение, либо удачное  
приложение теоретических принципов. 
Заинтересовать детский ум — это одно из основных  
положений нашей доктрины, и мы ничем не пренебрегаем, 
чтобы привить учащимся вкус, мы готовы сказать — 
страсть к учению». 
30 
«Итак решено, что если вместо краткого резюме, которое 
имеет целью ознакомить с сущностью вопроса, мы должны 
развернуть его историю, то мы не остановимся перед тем, 
чтобы изложить все методы, все попытки, все успехи, все 
заблуждения». 
«Изучение наук можно рассматривать с двух точек 
зрения. 
С одной стороны, хотят получить людей полезных и 
с опытом во всех областях деятельности цивилизованных 
народов. Речь идет о земледельцах, фабрикантах,  
коммерсантах, мореплавателях, офицерах, инженерах, врачах. 
С другой стороны, хотят иметь ученых, продолжающих 
отвлеченниые исследования, произведенные наиболее  
выдающимися умами, не дающих погибнуть результатам,  
накопленным в течение предыдущих веков; одним словом, 
речь идет о математиках и натуралистах, физиках и  
астрономах, об изобретателях, обладающих могучей логикой, 
мощным разумом, упорством для достижения цели, о  
неутомимых наблюдателях и глубоких вычислителях». 
«... путь, по которому идут, можно сравнить с тем, которым 
пользуются почти во всех военных школах. 
В этих школах намереваются подготовить офицеров, а 
поступают так, будто хотят подготовить исключительно 
генералов. 
Преподавая науки, хотят создать ученых, словно таких 
людей можно подготовить по своему желанию, и, стремясь 
к слишком многому, в результате не получают даже  
необходимого, т. е. полезных и скромных деятелей для наших 
современных общества. 
«...ученикам мы будем говорить: знайте немного, если 
вы не в силах знать больше, но знайте хорошо то, что вы 
выучили. Поэт, сказал: «Кто ограничивает себя, растет». 
Только небольшому числу людей дано знать много и 
хорошо. Все могут стремиться к этому, но только немногие 
достигают цели. Таким образом, учите хорошо только 
одно, если вы не можете запомнить больше; читайте всегда 
одну и ту же полезную книгу, вместо того, чтобы  
рассеивать свое внимание на нескольких. 
Ради себя и ради других отлично владейте своими  
специальными знаниями. 
Если вы будете первым в своей специальности, то будете 
более полезным для себя и для других, чем вы будете  
посредственно знать и науку, и литературу, технические 
ремесла и искусство». 
31 
«Образование прекращается только вместе с  
прекращением жизни. Ежедневно узнаешь что-то новое,  
отвлекаешься от своей обычной работы, делая что-то другое. 
И только глупец может считать, что с достижением 
определенного этапа нет больше ничего, что было бы  
полезно изучить». 
«...все дети любят физический труд. Они проворны, 
изобретательны, у них хорошее настроение до тех пор, 
пока школа не уничтожит в них большую часть этих  
драгоценных ростков». 
«...скука является самой опасной отравой. Она  
действует беспрестанно; она растет, овладевает человеком и 
влечет его к наибольшим излишествам». 
«...нужно полностью овладеть вниманием учащихся, 
направлять его, но при обязательном условии, чтобы  
постепенно возрастала сила суждения, без утомления и так, 
чтобы это не вызывало ни усталости, ни отвращения». 
«Дело не только в том, чтобы выучить — надо закрепить 
усвоенное. 
В этом состоит, по нашему мнению, наибольшая  
трудность преподавания...» 
«Мы предвидим серьезное возражение. 
«Вы рассуждаете, скажут нам; так, как будто уже  
созданы кадры воспитателей или преподавателей. Мы, как 
и вы, считаем, что хорошие преподаватели готовят  
хороших учащихся, но как обучить хороших  
преподавателей? Не будете ли вы вынуждены создавать образцовый 
педагогический институт, не придется ли еще ожидать 
целое поколение прежде, чем можно будет получить все 
то хорошее, что содержится в ваших предложениях?» 
«...Мы согласны, что это действительно очень серьезное 
возражение, и знаем, что необходимо обеспечить  
подготовку учителей и преподавдтелей...» 
«Преподаватель может, в крайнем случае, знать не более 
того, что он должен преподавать, при условии, что этими 
знаниями он обладает во всей их полноте, со всеми  
частностями, какие можно себе представить, и со всеми  
возможными непосредственными приложениями. 
Он может без особых неудобств ограничиться,  
например, тем, что будет знать арифметику, алгебру и основные 
понятия физики и химии, но он должен сочетать знание 
этих наук с полным их пониманием» 
32 
Никто в мире не должен знать об этих вещах больше, 
чем он, говорить о них лучше, сравнивать более  
старательно, писать более ясно...» 
«Преподаватель должен прежде всего любить 
свою профессию. Каждый как для своего личного 
счастья, так и для блага других людей должен любить 
свою профессию. Но преподаватель больше, чем кто бы 
то ни было, должен быть предан своей работе, считать ее 
целью всех своих усилий». 
«Мы страстно хотим приблизить ту пору, когда почти 
все люди науки, преданные своей родине, с  
воодушевлением займутся жизненно важным вопросом преподавания 
наук. 
Все упростится тогда в жизни нации, и наука станет - 
деятельным, настойчивым помощником, соучастником всех 
моральных и материальных достижений». 
А. Я. ХИНЧИН 
О ВОСПИТАТЕЛЬНОМ ЭФФЕКТЕ 
УРОКОВ МАТЕМАТИКИ 10 
Математика, в отличие от большинства других  
преподаваемых в школе дисциплин, имеет предметом своего изучения 
не непосредственно вещи, составляющие собою  
окружающий нас внешний мир, а количественные отношения и  
пространственные формы, свойственные этим вещам. Этой  
особенностью математической науки в первую очередь  
объясняются те хорошо известные методические трудности,  
которые неизбежно встают перед преподавателем математики 
и которых почти не знают преподаватели других наук: 
перед учителем математики стоит нелегкая задача —  
преодолеть в сознании учеников со стихийной неизбежностью 
возникающее представление о «сухости», формальном  
характере, оторванности от жизни и практики его науки. 
Об этом написано много ценного и полезного, и мы  
хорошо знаем, как справляются с этой задачей лучшие  
мастера нашей школы. 
Но этой же особенностью математической науки в  
значительной мере объясняется и специфика задач, встающих 
10 ХинчинА. Я- Педагогические статьи. М., Изд-во  
Академии пед.наук PCOCPg 1963, с. 128—160. 
33 
перед учителем математики, который хочет использовать 
преподавание своей науки в воспитательных целях. Ясно, 
что и здесь стоящая перед ним задача труднее, чем в случае 
большинства других наук. Ибо научная дисциплина,  
занятая изучением не самих вещей, а лишь отношений между 
ними и поэтому необходимо требующая подъема на  
некоторую ступень абстракции, — такая дисциплина, очевидно, 
лишь в редких случаях способна давать учителю повод к 
эффективному воздействию на формирование характера 
и мировоззрения учащихся, на регулирование их поведения. 
Этим, несомненно, объясняется то, что в исследованиях, 
посвященных вопросам воспитательной функции  
школьного обучения, об уроках математики обычно вовсе не  
говорится или говорится очень мало. 
То немногое, что написано по этому поводу, в основном 
не вызывает возражений. Дело сводится обычно к двум  
рычагам воспитательного воздействия: с одной стороны — 
говорится, что специфическая для математики логическая 
строгость и стройность умозаключений призвана  
воспитывать в учащихся общую логическую культуру мышления; 
с другой — указывается, что предметно-содержательное 
оснащение математцческих задач при надлежащем его 
выборе дает широкий простор для сообщения цифр и  
данных, способных значительно расширить кругозор  
учащихся, поднять их общий культурный урдвень и содействовать 
мировоззренческому и патриотическому моментам в их 
общем идейно-политическом воспитании. 
Это, бесспорно, верно, но, я думаю, далеко не все. 
Прежде всего здесь совершенно не затрагиваются  
важнейшие задачи морального воспитания, для которых, как мне 
представляется, уроки математики дают весьма  
ощутительные возможности. Далее, важная задача воспитания 
логической культуры мышления, которой обычно уделяется 
много внимания, тем не менее трактуется в большинстве 
случаев трафаретно, поверхностно и недостаточно  
расчленение; приводимые примеры часто не выходят за рамки 
вульгарного шаблона и поэтому очень мало эффективны. 
Наконец, воспитывающее воздействие данных,  
приводимых в «текстовых» задачах, хотя и должно всемерно быть 
использовано, но с математическим содержанием урока 
связано лишь весьма внешним образом. Ясно, что здесь 
воспитывающее влияние призвана оказывать не сама  
математика, не ее законы и ее стиль, а те привязанные к 
ней чисто внешним образом данные, которые обрамляют 
34 
сооою «текстовые» задачи и которые без всякого изменения 
математического содержания задачи могли бы быть  
заменены любыми другими аналогичными данными. Понятно 
поэтому, что этот рычаг воспитывающего воздействия,  
будучи важным и действенным, не может считаться в  
прямом смысле принадлежащим самой преподаваемой в школе 
математической науке. 
Все приведенные соображения показьГвают, что вопрос 
о воспитательном значении уроков математики у нас  
разработан еще далеко не достаточно. Предлагаемая статья 
ставит себе целью несколько осветить этот вопрос. Для 
этого я в дальнейшем кратко рассмотрю ряд моментов, 
которые, насколько я могу судить, при изыскании  
возможностей воспитательного влияния уроков математики до 
сих пор либо совсем оставлялись без внимания, либо  
рассматривались лишь поверхностно. 
КУЛЬТУРА МЫСЛИ 
Правильность мышления 
Роль и значение математики в воспитании навыков  
закономерного и безошибочного мышления в такой мере 
всеми признана, что нередко приходится встречаться с 
утверждениями, будто приучение к строгому в логическом 
отношении ходу мыслей есть первая и основная задача 
учителя математики, так что в сравнении с нею даже  
ознакомление учащихся с самим содержанием математической 
науки отодвигается на второй план (что, несомненно,  
следует признать уже вредным перегибом). Однако именно 
потому, что эта воспитательная функция уроков математики 
приобрела характер банальности, в этом направлении мы 
слышим много высказываний, приводимых по готовому 
трафарету, без достаточного обдумывания. В результате 
внимание сосредоточивается на небольшом числе  
привычных (а подчас и набивших оскомину), хотя и важных,- но 
по своему значению частных и узких вопросов, вроде,  
например, уже пресловутого различения между прямыми 
и обратными теоремами. Между тем оставляются в тени 
вопросы гораздо более принципиального значения. 
Я думаю, что основным моментом воспитательной  
функции математического образования — моментом, который 
в значительной степени обусловливает собою все осталь- 
35 
ное, — служит приучение учащихся к полноценности  
аргументации. 
В обыденной жизни, даже в «любительских» (не строго 
научных) принципиальных спорах, мы, защищая какое- 
либо утверждение, довольствуемся обычно одним-двумя 
аргументами, говорящими в его пользу. Противник может 
привести в ответ несколько аргументов, говорящих против 
нашего утверждения. Однако обычно ни та, ни другая 
аргументация не бывает исчерпывающей; противники  
продолжают изыскивать новые аргументы, каждый в пользу 
своей точки зрения, и спор продолжается. 
Примерно так же протекают и научные дискуссии 
в тех областях знания, которые не входят в число так  
называемых «точных» наук; конечно, аргументация здесь 
бывает уже, как правило, более полной, чем в обыденных 
спорах, но почти никогда не удается сделать ее  
исчерпывающей, не допускающей никаких возражений и тем  
самым ликвидирующей самую дискуссию. 
Иначе обстоит дело в математике. Здесь аргументация, 
не обладающая характером полной, абсолютной  
исчерпанности, оставляющая хотя бы малейшую возможность 
обоснованного возражения, беспощадно признается  
ошибочной и отбрасывается как лишенная какой бы то ни 
было силы. В математике нет и не может быть  
«наполовину доказанных» и «почти доказанных» утверждений: 
либо полноценность аргументации такова, что никакие 
споры о правильности доказываемого утверждения более 
невозможны, либо аргументация вообще отсутствует. 
Изучая математику, школьник впервые в своей жизни 
встречает столь высокую требовательность к полноценности 
аргументации. Вначале она удивляет, отталкивает, пугает 
его, кажется ему излишней, сверхмерной, педантичной. 
Но постепенно, день за днем, он к ней привыкает.  
Хороший учитель много может сделать для того, чтобы этот 
процесс протекал и быстрее и продуктивнее. Он приучит 
своих учеников к взаимной критике; когда один из них 
что-либо показывает или решает какую-либо задачу перед 
всем классом, все остальные должны напряженно искать 
возможных возражений и немедленно их высказывать. 
Ученик, который «отобьется» от всех таких возражений, 
заставит умолкнуть всех своих критиков, неизбежно  
испытает законную радость победы. Вместе с тем он ясно 
почувствует, что именно логическая полноценность  
аргументации была тем оружием, которое дало ему эту победу. 
36 
А раз почувствовав это, он неизбежно научится уважать 
это оружие, стараться, чтобы оно всегда было при нем. 
И, конечно, не только в математических, но и в любых 
других дискуссиях он все больше и настойчивее будет 
стремиться к полноценности аргументации. Каждый раз 
перед ним будет вставать задача — по возможности  
обезоружить своих противников, в полной мере используя весь 
запас аргументов, какие вообще мыслимы в данной ситуации. 
Этот воспитывающий процесс имеет решающее значение 
для логической культуры мышления, в особенности если 
учесть, что учащийся привыкает быть беспощадно  
требовательным к полноценности аргументации не только в споре, 
но и в своем одиноком мышлении. Процесс этот протекает 
повседневно на наших глазах у многих тысяч школьников. 
Он неизбежно возникает и идет своим путем без нашего 
специального вмешательства, но это не значит, конечно, 
что мы вправе предоставить его такому самотеку; в нашей 
власти сделать его и более быстрым, и более полным по 
богатству и прочности достижений; а раз мы можем, то мы, 
очевидно, и должны это делать: вопрос о том, какими  
приемами наиболее эффективно можно добиться этих целей, 
есть уже методическая задача, которую мы не можем здесь 
детально рассматривать. 
Общий принцип борьбы за полноценность аргументации 
получает в ходе интеллектуального развития учащегося 
целый ряд типичных по форме конкретных  
разновидностей, важнейшие из которых мы теперь перечислим. 
Борьба против незаконных обобщений. Натуралист, 
подметив наличие какого-либо свойства (признака) у ряда 
особей данного вида, с чистой научной совестью объявляет 
этот признак общим для всего рассматриваемого вида; 
и никто не упрекнет его за это — такого рода  
индуктивные заключения составляют собой один из основных  
методологических стержней естественных наук. Конечно, и в 
этих науках координирующая и осмысливающая  
теоретическая мысль возможна и необходима; но как исходным 
пунктом, так и решающей проверкой всякого заключения 
здесь всегда остаются наблюдение или опыт,  
осуществляемые над отдельными экземплярами. 
В математике дело обстоит принципиально иначе. Если 
мы проверили, что несколько десятков (или хотя бы и  
несколько миллионов) наудачу выбранных нами  
треугольников обладают каким-нибудь свойством, мы еще не вправе 
признать это свойство принадлежащим всем треугольни- 
37 
кам. Такое заключение было бы не до конца обоснованным, 
а в математической науке все, что не обосновано до  
конца, расценивается как абсолютно необоснованное. Только 
исчерпывающее общее доказательство может дать  
уверенность в том, что данный признак действительно является 
общим свойством всех треугольников. 
Чему же может и должна научить школьника та  
суровая критика по адресу не вполне обоснованных обобщений, 
с какою он встречается в математике? Конечно, он не 
должен стараться переносить такого рода требования на 
выводы других наук и тем более на практические  
жизненные ситуации. Требование абсолютной полноты индукции 
специфично для математического метода и совершенно 
невыполнимо ни в естественных науках, ни в  
практической жизни. Но привычка с критической тщательностью 
проверять законность всякого обобщения, привычка  
твердо помнить, что замеченное во многих случаях еще не  
обязано тем самым иметь место во всех случаях и что  
закономерности, установленные на основе (хотя бы и многих) 
единичных наблюдений и опытов, требуют поэтому все 
новой проверки — все эти важнейшие методологические 
навыки, необходимые в любой научной и практической  
деятельности, в значительной степени воспитываются и  
укрепляются вместе с повышением математической культуры. 
Это — процесс, который мы каждодневно -видим  
происходящим на наших глазах. 
Борьба против необоснованных аналогий. Заключения 
по аналогии служат обычным и законным приемом  
установления новых закономерностей как в эмпирических науках, 
так и в обыденной жизни. Если, допустим,  
естествоиспытатель помнит, что все встречавшиеся ему до сих пор 
виды, обладающие признаками А и Б, обладали также 
и признаком В, и если он нашел новый вид, у которого 
обнаружены признаки Л и Б, то он, естественно,  
заключит, что этот новый вид обладает также и признаком В. 
Такое заключение по аналогии значительно выигрывает 
в убедительности, если к чисто эмпирическим данным, 
описанным выше, присоединяются, как это часто  
бывает, какие-либо теоретические соображения,  
заставляющие предполагать, что совместное наличие признаков Л, 
Б и В является не случайным, а обоснованным теми 
или другими общими принципиальными соображениями. 
Но только в математике возможно — и вместе с 
тем совершенно необходимо — требовать, чтобы эти 
38 
принципиальные соображения были доведены до степени  
исчерпывающего доказательства. Либо мы со всей строгостью 
доказали, что из наличия признаков А и Б с неизбежностью 
вытекает и наличие признака В\ либо, если нам не удалось 
доказать этого с исчерпывающей полнотой, нам запрещается 
делать из наличия признаков А и Б какие бы то ни было  
выводы относительно признака В. Но в первом случае (т. е. когда 
доказана теорема «из А и Б следует В») простое  
применение этой общей теоремы к конкретным частным случаям 
уже вряд ли может быть названо «заключением по  
аналогии». Будет, таким образом, правильно сказать, что в 
математике заключения по аналогии категорически  
запрещены (что не должно, конечно, умалять огромного  
эвристического значения заключений по аналогии), в то время 
как в эмпирических науках и практической деятельности 
заключениям по аналогии принадлежит почетная роль 
одного из основных приемов вывода новых закономерностей. 
Поэтому снова встает вопрос о том, что же в этом  
отношении могут дать уроки математики для воспитания общей 
культуры мышления. И снова приходится ответить на 
это то же, что и прежде: математическая вышколенность 
ума, привыкшего к тому, что заключение по аналогии 
может служить лишь эвристическим приемом, который 
сам по себе еще не имеет доказательной силы, неизбежно 
приучает прошедшего эту школу человека и во всех  
других областях мышления относиться к такого рода  
заключениям с большой осторожностью, памятуя, что во всех 
таких случаях нельзя без основательной проверки  
считать полученное заключение твердо установленным.  
Каждый из нас испытал в свое время на себе воспитывающее 
влияние этой особенности математического мышления, и 
каждодневно мы наблюдаем, как влияние это содействует 
повышению мыслительной культуры наших  
воспитанников. Критическое отношение к заключениям по аналогии 
есть один из важнейших показателей, отличающих  
правильно воспитанное научное и практическое мышление от 
первобытного, обывательского; и занятия математикой 
всегда служат одним из основных средств воспитания 
атого важнейшего показателя. 
Борьба за полноту дизъюнкцийп. Когда математик 
доказывает какое-либо общее свойство всех треугольников, 
то иногда ему приходится проводить доказательство от- 
11 Дизъюнкция — разделение.— Б. Г* 
39 
дельно для косоугольных, прямоугольных и тупоугольных 
треугольников. Известно, как часто в таких случаях  
начинающие делают ошибки, в особенности в тех случаях,  
когда рассуждение сопровождается ссылкой на чертеж;  
чертится, например, косоугольный треугольник, и рассуждение 
опирается на добавочные построения, которые либо  
невозможны, либо теряют доказательную силу, если  
выбранный треугольник имеет тупой угол. В математике такое  
рассуждение признается ошибочным, так как здесь нарушено 
основное требование полноты дизъюнкции — не  
предусмотрены все возможные разновидности данной ситуации, 
одна из них выпала из поля зрения. 
В обыденных, не научных рассуждениях это требование 
нарушается на каждом шагу. Рассмотрев две-три наиболее 
часто встречающиеся или наиболее бросающиеся в глаза 
разновидности данной ситуации и убедившись, что в  
каждом из этих случаев мы неизбежно встречаемся с некоторы- 
рым событием Л, мы заключаем, что это событие А  
сопутствует данной ситуации во всех случаях, хотя на самом деле 
данная ситуация может иметь, кроме двух-трех изученных 
нами, еще десяток других разновидностей, и среди этих 
разновидностей, скинутых нами со счета, могут быть и 
такие, в которых наступление события А вовсе не  
обязательно. Мы говорим, например, что ученика Иванова  
вообще нельзя дисциплинировать, потому что на него  
испытанным образом не действуют ни ласка, ни угрозы. Мы  
забываем при этом, что лаской и угрозами не исчерпываются 
еще все разновидности приемов дисциплинирующего  
воздействия, что существует еще, например, метод спокойного 
убеждения и что, стало быть, наша дизъюнкция страдает 
неполнотой. Мы часто наблюдаем, как начинающий,  
рассмотрев при исследовании какого-нибудь уравнения  
случай, когда некоторый данный коэффициент положителен, 
а затем случай, когда этот коэффициент отрицателен, тем 
самым считает, что он провел исследование во всех  
случаях, забывая, что изучаемый коэффициент может оказаться 
равным нулю. Здесь также мы видим неполноту  
дизъюнкции, которая может привести и фактически приводит к 
тяжелым ошибкам в выводах. 
В противоположность тем двум требованиям, которые 
мы рассматривали выше, требование полноты дизъюнкции, 
учета всех возможных разновидностей изучаемой ситуации 
является необходимой принадлежностью не только  
математического, но и всякого правильного мышления. Аргу- 
40 
ментация, в которой не учтены все имеющиеся  
возможности, всегда оставляет место для законных возражений 
и потому не может быть признана полноценной.  
Военачальник, предпринимая какой-либо маневр, при учете его  
последствий должен предвидеть все возможные ответы врага; 
просмотр хотя бы одного из них может оказаться  
гибельным. 
Юридический кодекс в каждой статье обязательно  
должен охватывать все мыслимые разновидности данной  
ситуации, иначе он ставит судью перед необходимостью  
решать дела по своему произволу. 
Но нигде требование безукоризненной чистоты  
дизъюнкций не выставляется так явно и категорически, как в  
математике; и никто не обрушивается с такой быстротой и 
беспощадностью на замеченный просмотр в дизъюнкции, 
как вышколенный математик. Вот почему уроки  
математики должны воспитывать и действительно воспитывают в 
мышлении учащихся этот важнейший закон правильного 
рассуждения в несравненно большей мере, чем занятия 
другими предметами. 
Борьба за полноту и выдержанность классификации. 
Классифицирует не только ученый теоретик в своем  
кабинете; классификацией приходится очень часто заниматься 
и практическому работнику — инженеру, врачу, учителю, 
статистику, агроному. Общеизвестно, что невышколенный 
ум склонен допускать, производя классификацию, ряд 
типических ошибок; наиболее распространенными из 
таких ошибок являются нарушение полноты  
классификации и нарушение ее выдержанности, единопринципности. 
Нарушение полноты классификации состоит в том, что 
остаются понятия, не входящие ни в один из названных 
классов, и что, стало быть, названы не все классы. Простые 
примеры: на вопрос «какие ты знаешь растения?» — 
школьник отвечает «травы и деревья», забывая о  
кустарниках, лишайниках и многих других типах; войсковые 
части делятся на сухопутные, водные и воздушные  
(упускаются интендантские, части связи и многие другие);  
натуральные числа делятся на простые и составные (упускается 
число 1); вещественные числа делятся на положительные 
и отрицательные (упускается нуль). 
Требование полноты классификации формально  
аналогично рассмотренному нами выше требованию полноты 
дизъюнкции, но, конечно, отлично от него по содержанию. 
Там шла речь об обязательности охвата всех могущих 
41 
возникнуть ситуаций, здесь же — о необходимости  
перечисления всех разновидностей некоторого понятия. Но 
здесь, как и там, явно и неукоснительно требование  
полноты классификации провозглашается в математике  
преимущественно перед всеми другими науками, и поэтому 
уроки математики более всех других воспитывают в  
школьнике этот обязательный элемент правильного мышления. 
Требование выдержанности классификаций состоит в 
том, чтобы она проводилась по единому принципу, по  
единому признаку, Это требование, при строго правильном 
мышлении совершенно обязательное, очень часто  
нарушается не только в обывательских рассуждениях, но и в 
серьезной практике. Вот простые примеры такой  
невыдержанной классификации: суда делятся на весельные,  
парусные, моторные и военные; очевидно, классификация  
начата по принципу различных движущих сил, и последняя 
рубрика этот принцип нарушает; другой пример: обувь 
подразделяют на кожаную, брезентовую, резиновую и  
модельную — та же картина. Конечно, подобного рода  
перечисления не всегда претендуют на роль классификации, и 
в таких случаях соблюдение единого принципа не  
обязательно (например, объявление: завод приглашает на  
работу плотников, штукатуров, женщин и подростков). Но 
во всех случаях, когда такому перечислению приписывается 
классифицирующая функция, невыдержанность  
разделяющего принципа вызывает такую неотчетливость всей схемы, 
которая может привести и к теоретическим смешениям, и к 
практической путанице. Поэтому логически вышколенный 
ум всегда ощущает недостаток выдержанности  
классификации как существенный дефект рассуждения. И снова 
наиболее чувствительна к этому дефекту математическая 
наука, и поэтому именно на уроках математики школьник 
преимущественно развивает в себе эту потребность видеть 
всякую классификацию выдержанной, построенной на 
едином классифицирующем принципе. 
Я перечислил те моменты в борьбе за правильность 
мышления и полноценность аргументации, которые  
представляются мне наиболее важными. Как уже было сказано 
выше, я не могу входить в этой статье в обсуждение тех 
методических приемов, с помощью которых учитель  
математики может достигнуть наибольшего успеха в деле 
воспитания у своих учеников перечисленных мною  
моментов правильного мышления. Но я считаю необходимым 
сделать по этому вопросу одно методическое замечание об- 
42 
щего характера (для опытного учителя, впрочем,  
совершенно очевидное): все те требования правильного  
мышления, о которых шла речь выше, должны воспитываться в 
учащихся исподволь, от случая к случаю, без излишнего 
педалирования 12, не может быть и речи о том, чтобы 
посвящать специальный урок, например, борьбе с  
незаконными аналогиями; такая постановка дела может только 
безнадежно погубить весь ожидаемый эффект. Надо,  
напротив, всемерно избегать во всем этом деле общих  
рассуждений и обращать внимание учащихся на тот или другой 
логический момент исключительно на базе ярко  
убедительного конкретного математического материала. Потребность 
в логической полноценности аргументации воспитывается 
не постоянным надоедающим напоминанием о  
необходимости этой полноценности, а показом на конкретных  
примерах (поводы к которым дает почти каждый урок), как 
несоблюдение этого требования ведет к ошибкам и  
неувязкам. Надо не отвлеченно проповедовать полноценность 
аргументации, а приучить учащегося к тому, что каждый 
пробел в аргументации немедленно вызывает придирчивый 
вопрос со стороны учителя или — что много лучше — со 
стороны товарищей. 
Я не> буду говорить здесь о том, что следует  
использовать уроки математики для правильного понимания  
различия между прямым и обратным утверждениями, а также 
и ряда других аналогичных различений. С одной стороны, 
об этом так много уже писалось, что вряд ли я смог бы  
прибавить здесь что-нибудь новое. С другой стороны, мне 
представляется, что этого рода моменты, будучи, конечно, 
обязательными для логически правильного мышления, 
все же по своему частному, специальному характеру не 
имеют вне математики столь существенного значения, как 
те значительно более общие принципы, которые я  
перечислил выше. 
Стиль мышления 
Помимо специфических, особо строгих требований и  
логической правильности умозаключений, математика  
отличается от других преподаваемых в школе наук также и 
стилем своего мышления. Стиль этот, хотя и претерпевает 
J8 Педалирование — в данном случае означает подчеркивание, 
нажим,— Б. Л 
43 
на протяжении веков и даже десятилетий, довольно  
значительные изменения, все же имеет некоторые общие для 
всех эпох непреходящие черты, заметно отличающие его 
от стилей, принятых в других науках. 
Утвердившийся в той или другой науке стиль мышления 
не является, как можно было бы думать, только внешним 
и потому второстепенным фактором, имеющим лишь  
эстетическую ценность и не могущим поэтому существенно 
влиять на развитие данной науки. Напротив, стилем  
мышления в значительной степени определяется отчетливость 
теоретических связей, простота и ясность научных  
конструкций, наглядная конкретность понятий и многое другое, 
от чего в свою очередь зависят эффективность,  
плодотворность научных дискуссий и научного преподавания, а 
вместе с тем и темпы развития науки. Среди тех особых 
черт, которые присущи стилю математического мышления, 
имеется ряд таких, которым свойственно весьма общее и 
широкое значение; такая черта, еслГи она усваивается 
представителем какой-нибудь другой науки или  
практическим деятелем, оказывает нередко весьма существенные 
услуги как его собственному мышлению, так и усвоению 
его трудов учениками и последователями. Читая сочинения 
кого-либо из крупнейших классиков в другой научной 
области, математик подчас с некоторым удивлением  
восклицает: «Да ведь он мыслит совсем по-нашему!» —  
удивление происходит от того, что обычно в этой научной  
области принят совсем иной стиль мышления, имеющий очень 
мало общего с математическим. 
Но если усвоение некоторых черт математического 
мышления способно облагородить мыслительный стиль 
и в других областях знания и практической деятельности, 
сделать этот стиль более мощным и продуктивным орудием 
мысли, то очевидно, что не следует пренебрегать  
использованием уроков математики для приучения молодых умов 
к постепенному усвоению этих черт, к тому, чтобы эти 
черты стали прочными навыками их мышления — сначала 
в пределах математики, а потом и за ее пределами. Для 
того чтобы это сделать, надо в первую очередь постараться 
со всей тщательностью выявить те черты стиля  
математической мысли, о которых здесь идет речь. Это я теперь и 
постараюсь сделать. 
В основе каждого правильно построенного хода мыслей, 
независимо от предметного содержания его, лежит такая 
формально-логическая схема, вышколенным умом ощу- 
44 
щаемая как некий логический костяк, стройный и  
закономерный, обросший тем или другим конкретным  
содержанием. Независимо от стиля мышления, эта логическая 
схема должна быть закономерной, лишенной пробелов, — 
без этого рассуждение становится недоброкачественным и 
должно быть отвергнуто. 
Однако роль и положение этого логического скелета в 
данном ходе мыслей бывают весьма различны и  
существенным образом зависят именно от стиля мышления. В одних 
случаях логическая схема становится определяющим, 
руководящим моментом мышления, так что мыслящий все 
время имеет ее перед глазами и сообразно с нею выбирает 
и направляет и последовательные этапы рассуждения. 
В других, напротив, логический костяк остается  
затушеванным, мысль в гораздо большей степени направляется 
запросами конкретного содержания, роль логики сводится 
к последующему контролю, да и этот контроль в  
письменном или устном изложении часто только подразумевается 
и явно не проводится; логическая схема как целое остается 
вне поля зрения мыслящего. Разумеется, встречаются  
нередко и стили мышления, промежуточные между двумя 
указанными. 
Для математики характерно доведенное до предела  
доминирование логической схемы рассуждения; математик, 
потерявший, хотя бы временно, из вида эту схему, вообще 
лишается возможности научно мыслить. Эта своеобразная 
черта стиля математического мышления, в столь полной 
мере не встречающаяся ни в одной другой науке, имеет 
в себе много ценного. Очевидна, что она в максимальной 
степени позволяет следить за правильностью течения мысли 
и гарантирует от ошибок; с другой стороны, она  
заставляет мыслящего при каждой дизъюнкции иметь перед  
глазами всю совокупность имеющихся возможностей и  
обязывает его учесть каждую из них, не пропуская ни одной 
(такого рода пропуски легко возможны и фактически  
наблюдаются при других стилях мышления). Поэтому  
приобретенные на уроках математики стилистические навыки, 
связанные с описываемой чертой, имеют существенное 
значение для повышения общей культуры мышления  
учащихся. 
Очень интересным и ярким примером мышления в  
далекой от математики области, и тем не менее чрезвычайно 
насыщенного этой чертой, могут служить произведения 
Маркса. Читателя, который после изучения экономических 
45 
трудов ученых раскрывает «Капитал», с первых страниц 
поражает железная, непреклонная логика этих строк. 
Логическая схема с ее неумолимыми требованиями не 
только определяет ход мысли автора, но и настойчиво 
убеждает читателя, который не может уйти от ее  
направляющего влияния. Этот необычный для экономического  
сочинения стиль, почти приближающийся к математическому, 
неизменно вызывает в читателе ощущение прочности,  
надежности, предельной убедительности и в то же время 
много помогает ему в усвоении читаемого. 
Второю характерной чертой математического стиля 
мышления, о которой здесь должно быть упомянуто,  
является лаконизм, сознательное стремление всегда  
находить кратчайший, ведущий к данной цели логический путь, 
беспощадное отбрасывание всего, что не абсолютно  
необходимо для безупречной полноценности аргументации.  
Математическое сочинение хорошего стиля не терпит никакой 
«воды», никаких украшающих, ослабляющих логическое 
напряжение разглагольствований, отвлечений в сторону; 
предельная скупость, суровая строгость мысли и ее  
изложения составляют неотъемлемую черту математического 
мышления. Черта эта имеет большую ценность не только 
для математического, но и для любого другого серьезного 
рассуждения; лаконизм, стремление не допускать ничего 
излишнего помогает и самому мыслящему и его читателю 
или слушателю полностью сосредоточиться на данном 
ходе мыслей, не отвлекаясь побочными представлениями 
и не теряя непосредственного контакта с основной линией 
рассуждения. 
Корифеи научной мысли, как правило, мыслят и  
выражаются лаконично во всех областях знаний — даже  
тогда, когда мысль их создает и излагает принципиально 
новые идеи. Какое величественное впечатление  
производит, например, благородная скупость мысли и речи  
величайших творцов физики — Ньютона, Эйнштейна, Нильса 
Бора! Может быть, трудно найти более яркий пример того, 
какое глубокое воздействие может иметь на развитие  
науки именно стиль мышления ее творцов. 
В гораздо меньшей степени этот лаконизм присущ  
ораторским выступлениям. Здесь мы часто встречаем  
растянутость, излишнюю цветистость, пренебрежение  
прямотою логического пути в угоду украшающей образности 
(которой, конечно, нельзя отказать в присущей ей  
специфической силе воздействия). Однако и в этой области, 
46 
когда встает оратор, облекающий свою мысль в сжатую, 
скупую форму предельно кратких и неумолимо  
убедительных ходов, величественно жертвующий во имя этой 
железной логики всеми стилистическими «красотами»,  
всеми соблазнами красочной образности, мы видим, как  
внимание слушателей сразу подтягивается и напрягается, 
и чувствуем, что такая речь должна вызывать  
значительно большее доверие, а потому и оказывать большее  
воздействие, чем многие ярко-образные, оснащенные  
витиеватыми нагромождениями выступления, апеллирующие к 
чувству и воображению слушателей. 
Для математики лаконизм мысли является  
непререкаемым, канонизированным веками законом. Всякая  
попытка обременить изложение не обязательно нужными 
(пусть даже приятными и увлекательными для  
слушателей) картинами, отвлечениями, разглагольствованиями  
заранее ставится под законное подозрение и автоматически 
вызывает критическую настороженность. И поэтому  
именно уроки математики призваны дать учащимся,  
предпочтительно перед другими предметами, навыки лаконического, 
прямого, не знающего отвлечений, не обремененного  
никакими излишними моментами мышления. 
Далее, для стиля математического мышления  
характерна четкая расчлененность хода рассуждений. Если,  
например, при доказательстве какого-либо предложения мы 
должны рассмотреть четыре возможных случая, из  
которых каждый может разбиваться на то или другое число 
подслучаев, то в каждый момент рассуждения математик 
обязан отчетливо помнить, в каком случае и подслучае 
его мысль сейчас обретается и какие случаи и подслучаи 
ему еще остается рассмотреть. При всякого рода  
разветвленных перечислениях математик должен в каждый  
момент отдавать себе отчет в том, для какого вида родового 
понятия он перечисляет составляющие его видовые  
понятия. В обыденном, не научном мышлении мы весьма 
часто наблюдаем в таких случаях смешения и перескоки, 
приводящие к путанице и .ошибкам в рассуждении. Часто 
бывает, что человек начал перечислять виды одного  
какого-нибудь рода, а потом незаметно для слушателей 
(а часто — и для самого себя), пользуясь недостаточной 
логической отчетливостью рассуждения, перескочил в  
другой род и заканчивает заявлением, что теперь оба рода  
расклассифицированы; а слушатели или читатели не знают, 
47 
где пролегает граница между видами первого и второго 
рода. 
Для того чтобы сделать такого рода смешения и  
перескоки невозможными, математики издавна широко  
пользуются простыми внешними приемами нумерации  
понятий и суждений, иногда (но гораздо реже) применяемыми 
и в других науках. Те возможные случаи или те родовые 
понятия, которые надлежит рассмотреть в данном  
рассуждении, заранее перенумеровываются: внутри каждого  
такого случая те, подлежащие рассмотрению подслучаи, 
которые он содержит, также перенумеровываются (иногда, 
для различия, с помощью другой системы нумерации). 
Перед каждым абзацем, где начинается рассмотрение  
нового подслучая, ставится принятое для этого подслучая 
обозначение (например, ИЗ — это означает, что здесь 
начинается рассмотрение третьего подслучая второго  
случая, или описание третьего вида второго рода, если речь 
идет о классификации). И читатель знает, что до тех пор, 
покуда он не натолкнется на новую числовую рубрику, 
все излагаемое относится только к этому случаю и под- 
случаю. Само собой разумеется, что такая нумерация  
служит лишь внешним приемом, очень полезным, но отнюдь 
не обязательным, и что суть дела не в ней, а в той  
отчетливой расчлененности аргументации или классификации, 
которую она и стимулирует, и знаменует собою. 
Наконец, следует упомянуть еще об одной чисто  
внешней традиции математического стиля, могущей *при  
надлежащих условиях приобрести воспитательное значение, 
которым нельзя пренебрегать. Я имею в виду  
свойственную математике скрупулезную точность символики.  
Каждый математический символ имеет строго определенное 
значение: замена его другим символом или перестановка 
на другое место, как правило, влечет за собою искажение, 
а подчас и полное уничтожение смысла данного  
высказывания. Учащийся, непривыкший еще относиться с  
достаточной требовательностью к точности устной речи и  
письменного изложения, вначале может с некоторым  
легкомыслием отнестись к неуклонным и настойчивым  
приглашениям учителя математики — вести запись с абсолютной 
точностью; эти требования могут даже показаться ему 
педантичными и вызвать насмешку. Однако он очень 
быстро убедится на собственном опыте, что несоблюдение 
безукоризненной точности символической записи в  
математике влечет за собой немедленную расплату: он сам 
48 
теряет возможность понять смысл записанного, вынужден 
гадать, угадывает неверно и либо получает неправильный 
ответ, либо вообще лишает себя возможности решить  
задачу. В лучшем случае ему ценой значительных усилий 
удастся восстановить правильную запись и двигаться 
дальше, отправляясь от нее. 
Убедившись таким образом, что точность  
символической записи соответствует его собственным интересам, 
учащийся начинает следить за собою в этом направлении, 
и постепенно строгая правильность математической  
символики становится его привычкой. Но такого рода  
привычка, приобретенная в какой-либо одной сфере  
мышления, неизбежно приводит к воспитанию и общего стиля 
мышления учащегося; он начинает точнее выражаться 
и в устной речи, и в письменном изложении; в частности, 
он уделяет больше внимания правописанию,  
орфографические ошибки переживаются им с такой же остротой и  
таким же беспокойством, как математические. Мы неизменно 
наблюдаем, что ученики, научившиеся требовательно  
относиться к точности математической символики, легче и 
быстрее перестают делать орфографические ошибки. И я 
не знаю, возможно ли окончить школу, обладая требуемой 
для аттестата зрелости математической культурой и не 
научившись в то же время писать совершенно безошибочно. 
Заканчивая эту главу, посвященную вопросам  
воспитательного воздействия уроков математики на культуру 
мышления учащихся, я предвижу естественное и законное 
недоумение читателя по поводу того, что мною нигде даже 
не затронута проблема развития диалектического  
мышления. Я считаю себя обязанным дать по этому вопросу 
краткое разъяснение. 
Маркс и Энгельс с полным основанием утверждали, 
что математика не только дает для законов  
диалектического мышления богатейший иллюстративный материал, 
но систематически способствует развитию диалектических 
навыков мыслительного процесса. Однако, как  
неоднократно отмечалось основоположниками марксизма, в полной 
мере это может быть отнесено лишь к так называемой  
«высшей» математике, т. е. к математике переменных величин. 
Именно здесь мы приучаемся к математическому  
исследованию явлений природы и процессов техники в их  
живой изменчивости, а не статической неподвижности.  
Именно здесь величины исследуются в их взаимной  
зависимости (понятие функции), а не в отрыве друг от друга. Нигде 
49 
с такою наглядностью, как здесь, мы не видим в действии 
переход количества в качество, диалектический синтез 
первоначально антагонистических противоположностей и 
другие основные принципы диалектики. И это — одна из 
важнейших причин (впрочем, далеко не единственная), 
заставляющих нас признать абсолютно необходимым  
введение элементов высшей математики в курс средней  
школы 13. 
Но пока мы только боремся за это. Что же касается 
преподаваемой в школе «элементарной» математики, то 
и она, конечно, как всякая подлинная и живая наука, 
не лишена диалектических элементов. Но здесь они  
выступают разрозненно и с малой мощностью, и говорить о них 
в статье, посвященной лишь основным рычагам  
воспитательного воздействия уроков математики, я не решился. 
Впрочем, я имею в виду в ближайшем будущем составление 
другой статьи, специально посвященной вопросу о  
необходимости введения в школьное преподавание элементов 
высшей математики; в этой статье я надеюсь дать  
развернутую и убедительную картину того, каким мощным  
орудием воспитания навыков диалектического мышления  
могли бы стать уроки математики переменных величин14. 
МОРАЛЬНЫЕ МОМЕНТЫ 
И ВОСПИТАНИЕ ПАТРИОТИЗМА 
О роли и значении уроков математики в воспитании  
правильного и дисциплинированного мышления говорилось 
и писалось очень много. Напротив, о влиянии  
математических занятий на формирование характера и моральной 
личности учащегося не сказано почти ничего. Это вполне 
понятно: по абстрактности своего предмета  
математическая наука не может, конечно, давать учащемуся тех  
непосредственно впечатляющих, этически воздействующих 
и формирующих характер образов, картин, эмоций,  
какими обладают, скажем, уроки истории или литературы. 
Было бы, однако, весьма поверхностно делать отсюда 
тот вывод, что в деле формирования нравственной личности 
школьника уроки математики вообще должны быть  
скинуты со счетов. По моему многолетнему опыту работа 
*3 В новых программах это пожелание до некоторой степени 
учтено.— Б. Г. 
14 К сожалению, эта статья так и не была написана.— Б. Л 
60 
над усвоением математической науки неизбежно  
воспитывает — исподволь и весьма постепенно — в молодом 
человеке целый ряд черт, имеющих ярко моральную  
окраску и способных в дальнейшем стать важнейшими  
моментами в его нравственном облике. Сделать этот процесс 
более активным и результаты его более прочными —  
достойная задача для учителя. Но прежде всего надо  
тщательно разобраться в том, что это за черты и какие  
особенности математической работы способны их воспитывать. 
Честность и правдивость 
В обывательских тяжбах всякого рода каждая из  
спорящих сторон исходит, как правило, из желательного ей, 
выгодного для нее решения вопроса и с большей или  
меньшей изобретательностью изыскивает возможно более  
убедительную аргументацию для решения вопроса в свою 
пользу. В зависимости от эпохи, среды и содержания  
спора, стороны при этом апеллируют к тому или другому 
высшему авторитету — общечеловеческой морали,  
«естественному» праву, священному писанию, юридическому 
кодексу, действующим правилам внутреннего распорядка, 
а часто и к высказываниям отдельных авторитетных  
ученых или признанных политических руководителей. Все 
мы много раз наблюдали, с какою страстностью ведутся 
подобного рода споры и какой убежденностью дышит, по 
видимости, аргументация каждой из сторон; можно  
подумать, такой тяжущийся действительно обуреваем  
желанием найти и отстоять истинное, справедливое,  
отвечающее духу и букве призванного в качестве арбитра  
авторитетного источника, решение. 
Но хорошо известно, что подобную картину мы часто 
наблюдаем не в одних только обывательских тяжбах. 
В точности те же черты являет подчас и научная дискуссия. 
Выводы, с полной убежденностью сделанные одним  
ученым, с такой же убежденностью оспариваются другими; 
завязывается полемика, в которой каждая из сторон  
находит все новые и новые аргументы в пользу своей  
позиции — даже вновь поставленные опыты часто говорят 
каждому из спорящих как раз то, что ему желательно. 
В ходе полемики каждая из сторон не только стремится 
все более и более усиливать свою собственную позицию, 
но и стремится различными средствами дискредитировать 
позицию противной стороны, доходя иногда и до попыток 
61 
персональной дискредитации. И лишь сравнительно редко 
бывает, чтобы в такой затянувшейся полемике одна из 
спорящих сторон нашла честность и мужество признать 
свою позицию ошибочной. 
Субъективные основания такого рода явлений в жизни 
науки могут быть легко поняты: они ничем, к сожалению, 
не отличаются по своей неприглядности от субъективных 
оснований самых мелочных обывательских стычек. Что 
касается объективных оснований возможности подобного 
рода научных ситуаций, то и их нетрудно найти: в  
эмпирических науках всякая новая, еще не окончательно  
установленная закономерность фигурирует, по крайней мере 
временно, в качестве «рабочей гипотезы»; покуда вопрос 
не решен окончательно, имеются обычно как соображения 
(опытные и теоретические), говорящие в пользу этой  
гипотезы, так и такие, которые говорят против нее. Из двух 
ученых один может поставить своей задачей — собрать 
как можно больше аргументов, поддерживающих такую 
гипотезу, а другой — заняться собиранием фактов и  
соображений, способных вызвать к ней недоверие. Дело  
происходит как в уголовном процессе, где перед обвинителем 
и защитником ставятся задачи собрать, привести в  
порядок и изложить все аргументы, соответственно говорящие 
за и против виновности подсудимых. 
Само собой разумеется, что так поставленная научная 
дискуссия сама по себе не содержит еще ничего морально 
одиозного 1б. Собрать с возможною полнотою все  
имеющиеся аргументы за и против данной «рабочей гипотезы» — 
это во всех случаях приносило пользу прогрессу науки; 
нет, очевидно, ничего предосудительного и в том, что сбор 
аргументов за и против гипотезы выполняется двумя  
различными учеными (или группами ученых), если только 
обе стороны подходят к своей задаче добросовестно,  
руководясь исключительно желанием способствовать  
отысканию объективной истины. Моральный одиум, этическое 
неблагополучие начинается там, где выводы ученого  
перестают руководствоваться интересами объективной истины, 
а становятся — сознательно, полусознательно или  
бессознательно — на службу его личным интересам — его 
упрямству, его честолюбию, его корыстолюбию, когда 
аргументация приводится с пристрастием, притягивается 
15 Одиозный— нежелательный, неприемлемый, вызывающий 
к себе отрицательное отношение.— Б. Л 
62 
за волосы, необъективно акцентируется, точь-в-точь как 
в обывательских дрязгах. Такая деградация научного 
спора в иных случаях ложится мрачным пятном даже, на 
крупнейших представителей научной мысли. 
Однако только математическая наука полностью от 
всего этого избавлена. Она не знает рабочих гипотез — 
предложений, истинность которых может подлежать  
дискуссии. Пока предложение не доказано, оно вообще не 
входит в сокровищницу науки, никому не придет в голову 
его отстаивать; если же оно доказано, то истинность его 
никак не может быть подвергнута сомнению: оно является 
абсолютно обязательным. Никаких промежуточных  
ситуаций математика не знает. Полемизировать, например, 
в защиту неполноценного доказательства в математике 
может только неуч, шарлатан или душевнобольной (все 
три категории действительно от времени до времени  
встречаются, достаточно вспомнить так называемых «фермати- 
стов», рыцарей квадратуры круга и трисекции угла); 
но такой «защитник» немедленно, единогласно и  
беспощадно разоблачается научным миром. Никакая  
аргументация с пристрастием или с тенденцией, никакое  
«притягивание за волосы» ни при каких обстоятельствах не может 
в математике иметь успеха. Разумеется, это относится  
только к содержанию самой математической науки; в вопросах 
логического или философского обоснования математики 
дискуссии возможны и даже неизбежны; возможны и 
споры персонального характера, связанные с развитием 
математики (например, по вопросам приоритета)16. 
Каждый математик рано привыкает к тому, что в его 
науке всякая попытка по тем или иным мотивам  
действовать тенденциозно, заранее склоняясь к тому или другому 
решению вопроса и прислушиваясь только к аргументам, 
говорящим в пользу избранного решения,— всякая такая 
попытка заведомо обречена на неудачу, и ничего, кроме 
разочарования, принести не может. Такое положение, 
при котором неправильная или не до конца правильная 
аргументация могла бы оказаться выгодной для  
аргументирующего, здесь просто принципиально невозможна.  
Поэтому математик быстро привыкает к тому, что в его науке 
16 Конечно, и математики способны на неприглядные  
житейские споры, а также на споры в области педагогики, во время  
которых некоторые из них забывают о принципе полноценности  
доказательств.— Б. Г. 
53 
выгодна только правильная, объективная, лишенная вся* 
кой тенденциозности аргументация, что успех может  
принести только непредубежденное, беспристрастное  
напряжение мысли. И независимо от своего общего морального 
уровня, он в своей научной работе всегда руководствуется 
исключительно соображениями объективной истинности. 
Но эту черту, естественно развивающуюся у  
математика-специалиста, в известной степени воспитывает в себе, 
занимаясь математикой, и каждый неспециалист, в  
частности, каждый школьник. Ему хорошо известно, что  
втереть очки учителю математики невозможно, что никакой 
апломб и никакое красноречие не помогут ему выдать 
незнание за знание, неполноценную аргументацию 5а 
полноценную. И как бы лжив он ни был в других  
отношениях, в математике он остережется отстаивать неверное 
утверждение или неправильное доказательство. 
Но и здесь, как это часто бывает, моральные навыки, 
приобретенные в какой-либо одной области, в известной 
мере переносятся и на другие сферы мышления и  
практической деятельности. Теоретическая честность, ставшая для 
математика непреложным законом, его научного  
мышления и профессиональной (в частности, педагогической) 
деятельности, довлеет над ним во всех его жизненных 
функциях — от абстрактных рассуждений до  
практического поведения. 
Я должен признаться, что органически не способен 
отстаивать какое-либо утверждение (хотя бы и обыденно- 
практического содержания), если я не располагаю не 
допускающим никакого возражения его доказательством. 
Профессиональная привычка к абсолютной объективности 
аргументации не позволяет мне, как это делают многие 
другие, яростно, во что бы то ни стало отстаивать  
выгодное мне решение. Таким образом, черта, о которой я  
сейчас говорю, может иногда и повредить своему носителю; 
тем не менее я дорожу ею и горжусь, что она у меня есть; 
радуясь и тогда, когда вижу ее в других, потому что  
придаю ей высокую моральную ценность. 
Я всегда интересовался этой чертою и много раз  
наблюдал, как она развивается в людях под влиянием  
серьезного научного общения, в частности, под воздействием 
уроков математики. Это очень радостная и морально  
возвышающая картина, когда человек постепенно  
преодолевает в себе отвратительную мещанскую- привычку —  
подчинять законы мышления своим личным, мелким, корыст- 
54 
ным интересам, теоретически защищать все то, и только 
то, что ему практически выгодно; когда он приучается 
уважать объективную правильность аргументации как 
высшую духовную и культурную ценность и все чаще 
и со все более легким сердцем жертвовать ради нее  
своими личными интересами. Доведенная до предела, эта 
черта составляет собою ничто иное, как честность и  
правдивость — одно из лучших украшений нравственной  
личности человека. 
Настойчивость и мужество 
Добросовестная и серьезная работа над приобретением 
и укреплением знаний в любой научной области требует 
систематического напряжения умственных усилий,  
настойчивости в преодолении трудностей, мужественной 
встречи неудач; поэтому такая работа при правильном 
руководстве неизбежно воспитывает у учащегося  
соответственные черты характера — трудолюбие, усидчивость, 
упорство в преследовании намеченной цели, умение не 
останавливаться перед трудностями и не впадать в уныние 
при неудачах. Непосредственно ясно, какое решающее 
значение имеют все эти черты для развития морально, и 
общественно полноценной человеческой личности и с  
каким вниманием поэтому должен учитель следить за  
максимальным использованием своих уроков в целях  
воспитательного воздействия в указанном направлении. Те  
возможности, которыми для этого располагают предметы 
школьного обучения, весьма разнообразны и  
многочисленны, и нет такого предмета, в специфических чертах 
которого не было бы заложено особых, именно этому  
предмету свойственных движущих рычагов такого  
воспитательного воздействия. Наша задача здесь, естественно, должна 
состоять в указании тех черт математики как школьного 
предмета, которые, отличая ее от других предметов  
школьного преподавания, наиболее способствуют развитию в 
учащихся разумной настойчивости и сознательного  
мужества — этих неоценимых качеств будущего борца. 
Прежде всего я хочу здесь отметить четкую  
определенность поставленной цели, желаемого и требуемого  
результата каждого математического задания. Если заданием 
служит сочинение исторического или литературного  
содержания, то нельзя указать момента, когда такое задание 
55 
дефинитивно17 закончено выполнением,— возможности  
дополнения и совершенствования, систематические  
улучшения всякого рода здесь почти безграничны; с другой  
стороны, учащийся не чувствует себя здесь достаточно  
компетентным для авторитетной оценки своей работы: то, что 
ему представляется в сочинении вполне удачным, может 
встретить совсем иную оценку со стороны учителя. Вся 
эта по существу данного задания неизбежная  
неопределенность, расплывчатость в оценке законченности и качества 
проделанной работы должна, несомненно, оказывать  
некоторое расслабляющее влияние на волевое напряжение 
еще мало вышколенного молодого ума. В математике дело 
обстоит иначе. Если заданием служит решение задачи 
или доказательство теоремы, то тем самым указывается 
с полной определенностью и тот момент, когда задание 
может считаться окончательно выполненным: когда решена 
задача или доказана теорема; все остальное — изложение 
найденного решения, правильность и аккуратность  
записи и т. п. — имеет и в глазах учителя, и в глазах  
ученика лишь второстепенное, не решающее значение.  
Равным образом и качество работы здесь оценивается с  
однозначной определенностью: задача должна быть решена 
верно, теорема должна быть доказана правильно.  
Проверить отсутствие логических ошибок в своем рассуждении 
ученик может и должен уметь сам; в случае задачи он 
знает даже определенные приемы проверки решения. 
Легко понять, какое стимулирующее влияние на упорство, 
настойчивость в достижении цели может оказать и  
действительно оказывает эта четкая определенность  
показателей результата. Победа здесь так же непосредственно 
ощутительна, как в шахматной партии или спортивном 
состязании, и сам учащийся может с такой же уверенностью 
зафиксировать и оценить свое достижение, как и его  
авторитетный учитель. 
Вторая, значительно более глубокая и важная черта 
математических заданий, которую я хочу здесь отметить, 
состоит в присущем в значительном большинстве случаев 
творческом характере. В то время как в большинстве  
других областей знания выполнение задания, за немногими 
исключениями, требует от учащихся лишь определенных 
знаний и навыков, в лучшем случае еще умения стройно 
и стилистически правильно излагать эти знания,— реше- 
17 Дефинитивно — с полной определенностью.— Б. Г. 
56 
ние математической задачи, как правило, предполагает 
изобретение специального, ведущего к поставленной цели 
рассуждения, и тем самым становится — пусть весьма 
скромным — творческим актом. Именно этот творческий, 
исследовательский характер математических заданий  
более, чем что-либо другое, влечет к себе молодые силы 
растущего и крепнущего интеллекта учащегося. Тот, 
кто раз изведал благородную радость творческого  
достижения, никогда уже не пожалеет усилий, чтобы вновь 
ее испытать. Никакие трудности его не остановят, сила 
его порыва и устремления, его усидчивость и выдержка 
в преодолении препятствий будут крепнуть с каждым  
новым достижением, а неудачи, ошибки, временные  
крушения и поражения он научится встречать как подобает 
истинному борцу — не опуская перед ними руки, а черпая 
в них источник и стимул для все новых и новых  
напряжений мысли и воли. 
Воспитание патриотизма 
Задача использования уроков математики для воспитания 
и укрепления в учащихся прочного чувства гордости 
своей Родиной и любви к ней имеет в себе специфическую 
трудность, очевидная причина которой заложена в  
абстрактном характере математической науки. Надо прямо  
сказать, что непосредственно, своим собственным материалом 
и содержанием математика в силу этой причины вообще 
не может служить орудием пропаганды чего-либо столь 
конкретного, как красота и величие родной страны. Здесь 
она с естественной скромностью вынуждена уступить 
место другим наукам. 
Однако на уроках математики ученик вовсе не все 
время сосредоточивается на ее абстрактной сущности;  
абстрактные схемы математики непрестанно, почти на  
каждом уроке оснащаются, дополняются и иллюстрируются 
весьма различным конкретным содержанием; сюда входит 
содержательный материал «текстовых» задач,  
исторические сведения, различного рода приложения и т. д. При 
этом во многих случаях выбор конкретного оснащения в 
весьма широких пределах может быть варьирован и таким 
образом в значительной степени ставится на усмотрение 
преподающего. Очевидно, такой произвол может быть  
широко использован учителем для фиксирования внимания 
учащихся на фактах и цифрах, поддерживающих и уваже- 
57 
ние и любовь к Отечеству. У нас неоднократно писалось 
уже о подборе патриотически направленного материала 
текстовых задач. Против этого приема ничего нельзя  
возразить; надо только тщательно продумать выбираемый 
материал, чтобы избежать опошления, вульгаризации самой 
патриотической идеи, как это бывает, когда конкретное 
содержание задачи малоестественно, «притянуто за  
волосы», или когда задача, сообщая достаточно интересные 
цифры и факты, ставит по поводу них такой вопрос,  
который явно не имеет ни непосредственного интереса, ни  
какого-либо практического зна'чения. Вместе с тем надо, 
конечно, отчетливо представлять себе, что весь этот прием 
является чисто внешним, для развития патриотических 
чувств здесь используются уроки математики, но никак 
не самая математика. 
Значительно теснее связан с самой математической 
наукой прием, состоящий в придании патриотической  
направленности целому ряду исторических сведений. Этот 
прием помимо впечатляющей силы воздействия, особенно 
ценен еще тем, что он значительно повышает интерес  
учащихся к истории математической науки, а во многих  
случаях дает повод и возможность эффективным образам 
ознакомить учащихся с математическими фактами,  
выходящими за пределы официальной программы и счастливым 
образом ее дополняющими. Так как по этому поводу у 
нас почти ничего не писалось, то я остановлюсь на этом 
несколько подробнее. 
История русской и советской математики богата  
фактами, знакомство с которыми, в особенности на фоне  
правильной исторической перспективы, способно возбуждать в нас 
законную радостную гордость. И среди этих фактов есть 
немало таких, понимание которых доступно учащимся 
средней школы в достаточной мере для того, чтобы они 
могли оценить их принципиальное или практическое 
значение. Нужно только, чтобы сам учитель был хорошо 
осведомлен как о самих этих фактах, так и об их роли 
и месте в науке, а также о той научно-исторической  
обстановке, в которой они возникали и развивались. Нужно, 
кроме того, уметь рассказать учащимся об этих фактах 
так, чтобы возбудить их живой интерес и извлечь  
максимальный эффект как для их математического развития, 
так и для воспитания в них здорового чувства  
национальной гордости. 
58 
Хорошо известно, что для всего этого очень продуктивно 
могут быть использованы научные идеи нашего великого 
соотечественника Н. И. Лобачевского и научная судьба 
его идей. В основе своей великий геометрический замысел 
Лобачевского вполне доступен школьникам старших  
классов, а проведенная с надлежащим тактом беседа о нем 
может много содействовать, с одной стороны, пониманию 
основной для современной математики идеи  
аксиоматического мышления, а с другой — глубокому уважению как 
к научному гению Лобачевского, так и к его замечательной 
теоретической стойкости — великой силе убеждения,  
позволившей ему творить в одиночестве, без общественного 
признания, в научно враждебной атмосфере. 
Значительно менее известны у нас творения другого 
нашего великого ученого — П. Л. Чебышева. А между 
тем научный облик его не менее импозантен, чем фигура 
Лобачевского. И кое-что о нем с большой и многосторонней 
пользой может быть рассказано и школьникам. Чебышев 
принадлежал к числу тех немногих ученых самого  
высокого ранга, которые на протяжении своей жизни работают 
в довольно многих, часто весьма удаленных друг от друга 
областях математики, в каждой из этих областей  
прокладывая совершенно новые пути, по которым затем, в течение 
многих десятилетий идут их последователи. Великий дух 
новаторства был присущ Чебышеву не в меньшей степени, 
чем Лобачевскому. В теории чисел, теории вероятностей, 
теории механизмов и теории аппроксимации функций он 
создал мощные новые методы и сделался родоначальником 
большого числа научных школ в России и за границей. 
Замечательные идеи его далеко не исчерпаны и до  
настоящего времени. 
Для учащихся средней школы особенно доступны и  
поучительны достижения Чебышева в теории чисел. Теорему 
Евклида о существовании бесконечного множества  
простых чисел знают все. Очень полезно выписать с учащимися 
таблицу простых чисел хотя бы до 100 и обратить их  
внимание на видимое отсутствие закономерности в  
расположении этих чисел. Затем рассказать о том, как задача о  
закономерностях в чередовании простых чисел была и  
остается одной из центральных проблем арифметики. Стоит  
привести (без доказательства) вполне понятный школьникам 
и способный вызвать в них интерес результат Эйлера 
59 
я (ri)!n -> 0 (п -+• оо)18. В самых общих чертах можно 
затем коснуться асимптотических результатов Чебышева, 
обязательно давая историческую картину тех значительных 
усилий, которые до Чебышева были посвящены этой  
задаче. Конкретно же очень стоит остановиться на  
элементарном постулате Бертрана, проверить его на ряде примеров 
и тем возбудить интерес к нему со стороны учащихся х 9. 
Позднее можно разобрать и какое-либо из его  
элементарных доказательств хотя бы в порядке кружковой работы. 
Очень советую обратить внимание учащихся на  
следующий замечательный исторический факт. Арифметика и 
геометрия — два старейших и важнейших раздела  
математической науки, и в обоих в течение ряда столетий 
наука в значительной степени питалась творениями  
Евклида; центральные проблемы этих двух основных ветвей 
математики — теория параллельных в геометрии и задача 
о распределении простых чисел в арифметике — в течение 
многих веков не поддавались сколько-нибудь заметно 
многочисленным усилиям целых поколений ученых. И вот 
в XIX столетии обе эти проблемы были сдвинуты,  
наконец, с мертвой точки. В геометрии это сделал русский 
математик Лобачевский, в арифметике — русский  
математик Чебышев. Оба они проложили, каждый в своей 
области, совершенно новые пути — пути, по которым  
наука успешно развивается до настоящего времени. Нет 
сомнения в том, что эти великие исторические скачки — 
Евклид — Лобачевский и Евклид — Чебышев — должны 
импонировать молодым умам, которые в известной мере 
уже способны оценить их значение. 
Заинтересовав учащихся вопросами распределения 
простых чисел, учитель имеет совершенно естественный 
повод рассказать им о знаменитой гипотезе Гольдбаха. 
Очень стоит проверить ее в классе в пределах хотя бы  
чисел первой сотни. Затем, конечно, без всяких  
доказательств, сообщить о блестящих достижениях советского 
академика И. М. Виноградова (его основной результат 
в направлении проблемы Гольдбаха, разумеется, вполне 
доступен учащимся по своему содержанию). 
18 Здесь я (п) обозначает число простых чисел между 2 и п.— 
Б. Г. 
19 Постулат Бертрана состоит в следующем: между числами 
а и 2а (а>>2) лежит хотя бы одно простое число. Эта гипотеза была 
доказана Чебышевым в качестве следствия из его общих  
результатов.— Б. Г. 
60 
ЗАКЛЮЧЕНИЕ 
Я сознательно полностью оставил в стороне важнейший 
для темы настоящей статьи вопрос о значении уроков 
математики для формирования научного мировоззрения 
учащихся. Я сделал это по той же причине, по какой в 
свое время (в первом разделе) отказался от рассмотрения 
вопроса об использовании уроков математики для  
воспитания навыков диалектического мышления: ознакомление 
с идеями и методами математической науки имеет  
фундаментальное мировоззренческое значение, но львиная доля 
воспитательного эффекта в этом направлении  
принадлежит математике переменных величин, так называемой 
высшей математике с которой, по выражению Энгельса, 
в математику входит, диалектика и которая, к сожалению, 
все еще почти целиком остается за бортом наших школьных 
программ. Поистине не стоит при этих условиях говорить 
о влиянии на мировоззрение формирования той  
диалектически худосочной части математической науки, которая 
в настоящее время находит себе приют в нашем школьном 
преподавании. О тех же мощных сдвигах, направляющих 
и формирующих мировоззрение, которое может произвести 
и фактически часто производит изучение высшей  
математики, можно сказать очень много, и я собираюсь говорить 
о них в моей будущей статье, о которой я уже упоминал 
на с. 50. 
В этой статье я не касался методических вопросов. 
Я говорил только о том, какие особенности математической 
науки и для воспитания каких именно качеств интеллекта 
или моральной личности учащегося могут и должны быть 
использованы, но нигде не касался вопроса о том, как это 
может быть сделано. Мой опыт говорит мне, что это  
обстоятельство может вызвать недовольство в некоторых кругах 
читателей; вероятно, я не гарантирован от набивших оско- 
61 
мину упреков в том, что моя статья «ничего не дает  
учителю», и советов «повернуться лицом», и т. д. 
Поэтому я считаю необходимым сделать по этому 
поводу следующее краткое разъяснение. 
1. Я считаю, что составление сколько-нибудь  
детальных методических указаний по рассматриваемым мною 
в настоящей статье вопросам действительно совершенно 
излишне. Сколько-нибудь заметный воспитательный  
эффект уроки математики (как и всякой другой науки) могут 
дать только при том условии, что учитель, во-первых, 
достаточно хорошо знает свою науку, ее методологию и ее 
историю, во-вторых, имеет достаточный педагогический 
такт и опыт и, в-третьих, сам обладает в достаточной мере 
всеми теми качествами, которые он собирается  
воспитывать в своих учениках. Учителю, который сам не умеет 
мыслить абсолютно отчетливо, никакие методические  
шпаргалки не помогут воспитать ясность мысли в учащихся; 
или еще: какая методика поможет учителю сделать  
учеников горячими патриотами, если сам он любит свою Родину 
вяло и с прохладцей? 
И, напротив, если учитель стоит на высоте своей задачи, 
если он в полной мере обладает всеми перечисленными 
выше качествами, то никакие методические разработки 
по воспитательным вопросам ему заведомо не нужны — 
в каждом отдельном случае он с легкостью и  
непринужденностью сам найдет наиболее эффективный путь к  
поставленной цели. Навязывание ему определенной конкретной 
методики было бы для такого учителя только помехой 
в его работе. 
2. Если, таким образом, я считаю составление  
детальных методических указаний по затронутым мною вопросам 
практически бесцельным, то, может быть, было бы полезно 
все же дать по ним ряд общих методических советов; я  
думаю, что было бы хорошо, если бы моя статья побудила 
кого-либо из наших лучших учителей-методистов  
высказаться по этому вопросу и сделать достоянием младших 
товарищей некоторые выводы из своего опыта. В этом они 
во всяком случае компетентнее меня, и здесь я со всей  
необходимой скромностью должен уступить им слово. 
Наконец, я хочу попытаться заранее оградить себя еще 
от одного рода упреков, которые я предвижу и которые 
обычно бывают основаны на недоразумении. Так как 
я-говорил о воспитательном эффекте уроков математики, 
то мне, естественно, пришлось перечислить одну задругой 
62 
именно те черты математической науки, которые в  
воспитательном отношении дают ей то или иное преимущество 
перед другими дисциплинами. В этом ведь и состояла моя 
вадача. Но когда поступаешь таким образом, то у  
недостаточно вдумчивого читателя создается впечатление,  
будто бы. ты поставил своей целью превознесение математики 
над всеми другими науками и вся твоя статья сплошь 
утверждает, что единственная подлинная наука есть именно 
математика, все же другие дисциплины страдают теми 
или иными изъянами и науками могут быть названы 
лишь с известной оговоркой. Уважаемые коллеги —  
представители других наук — начинают чувствовать себя 
несправедливо обиженными и подвергают твою работу 
яростной критике, доказывая, что другие науки ничуть 
не хуже и что всеми преимуществами, которыми в моем 
представлении монопольно обладает математика, на самом 
деле в такой же мере наделены и все прочие дисциплины. 
С первым пунктом этого заявления я целиком и  
полностью согласен — другие науки действительно ничем 
не хуже математики; более того, представители этих  
других наук обычно вызывают во мне глубочайшее уважение 
тем, что творят великие ценности в таких областях, в  
которых творчество, на мой взгляд, неимоверно трудно — 
гораздо труднее, чем в математике. Но ведь в моей статье 
я нигде ни разу не называю математику лучшей из наук. 
Напротив, я несколько раз со всею скромностью  
подчеркиваю, что как орудие воспитания математика лрежде всего 
отмечена такою особенностью, значение которой для  
данной цели очевидным образом отрицательно — она  
абстрактна; предметом ее служат не сами вещи и явления 
реального мира, а лишь абстрагированные от них  
количественные отношения и пространственные формы. Это 
обстоятельство, как я несколько раз подчеркиваю в своей 
статье, делает для математики воспитательную задачу 
значительно труднее, чем для других школьных  
дисциплин. Но зато математика в некоторых других отношениях 
отмечена такими чертами, которые создают ей  
воспитательные возможности более значительные, чем у других 
дисциплин. И то, что я в своей работе, соответственно 
моей задаче, сосредоточиваю внимание читателя именно 
на этих чертах,— никак не может, конечно, означать 
какого-либо гипертрофирования математики,  
превознесения ее выше всех других наук. 
Я отдаю математической науке лишь то, что ей принад- 
63 
лежит по праву, с полной откровенностью признавая и те 
ее черты, которые в отношении данной цели составляют 
ее слабость. Но за те преимущества, которые ей мною 
приписаны, я уже действительно готов драться до конца, 
И если уважаемые коллеги пожелают утверждать, что та 
или другая конкретная черта, по моему утверждению, 
монопольно или преимущественно присущая математике, 
на самом деле является в такой же мере достоянием и 
всех других научных дисциплин, то здесь я где угодно 
и когда угодно готов держать ответ в полной уверенности, 
что смогу отстоять правильность моего утверждения перед 
любым компетентным и беспристрастным судом. 
МАТЕМАТИКА КАК ПРОФЕССИЯ 
Сборник статей 
Составитель — ГНЕДЕНКО Б. В. 
Главный отраслевой редактор Л. А. Е р л ы к и н 
Редактор Г. Г. Карвовский 
Мл. редактор Т. Г. Иншакова 
Обложка художника Л. П. Ромасенко 
Х/дож. редактор М. А. Бабичева 
Техн. редактор А. М. Красавина 
Корректор В. И. Гуляева 
ИБ № 2800 
Сдано в набор 10.04.80 г. Подписано к печати 21.05.80 г. T-10165. 
Формат бумаги 84ХЮ87з2. Бумага № 3. Гарнитура литературная. 
Печать высокая. Усл. печ. л. 3,36 Уч.-изд. л. 3,57. Тираж 35 720. 
Заказ № 818. Цена 11 коп. Издательство «Знание». 101835, ГСП 
Москва, Центр, проезд Серова, д. -4. Индекс заказа 804306. 
Чеховский полиграфический комбинат Союзполиграфпрома 
Государственного комитета СССР 
по делам издательств, полиграфии и книжной торговли 
г. Чехов Московской области 
11 коп. 
Индекс 70096

	
	\newpage
	\tableofcontents
	
	\thispagestyle{empty} % 
	
	\newpage
	
	\setcounter{secnumdepth}{0}
	
	\phantomsection
	
		\section*{Описание}
	
	{\bf Название:} Библия людоедов
	
{\bf Автор:} Александров Георгий Федорович

{\bf Редактор:} Л. Лазаревский
	
{\bf Издательство:} Москва : Государственное издательство политической литературы, 1942. — 62, [2] с.
	
		{\bf Аннотация:} Жалкие выродки истории, как метко назвал немецких фашистов Горький, завсегдатаи кабаков и участники почти всех преступлений, совершенных за последние годы в Германии, разбойничьи захватив власть и назвав себя всем на посмешище «национал-социалистами», вот уже несколько лет организуют одну кровавую бойню за другой.
		


		\thispagestyle{empty} % выключаем отображение номера для этой страницы


\newpage

\setcounter{secnumdepth}{0}  

\phantomsection	

\ \\
\ \\
\ \\
\ \\

\hangindent=2cm \hangafter=0  	\textbf{ИНФОРМАЦИЯ ОБ ОЦИФРОВКЕ}

\hangindent=2cm \hangafter=0  	\textbf{Источник}:  \href{https://rutracker.org/forum/viewtopic.php?t=5855912}{https://rutracker.org/forum/viewtopic.php?t=5855912} 

\hangindent=2cm \hangafter=0  	\textbf{Версия}: 1.0 - от 02 октября 2026

\hangindent=2cm \hangafter=0  	\textbf{Оцифровка}: lord199 (\href{mailto:lord199@mail.ru}{lord199@mail.ru})



\hangindent=2cm \hangafter=0  Если вы нашли какие-то опечатки и другие неточности, просьба связаться по электронному адресу.




\thispagestyle{empty} % выключаем отображение номера для этой страницы
	
\end{document}


